\documentclass[conference,a4paper]{IEEEtran}
\IEEEoverridecommandlockouts

\usepackage[utf8]{inputenc}
\usepackage[spanish]{babel}
\addto\captionsspanish{\renewcommand{\tablename}{Tabla}}
\usepackage{cite}
\usepackage{amsmath,amssymb,amsfonts}
\usepackage{graphicx}
\usepackage{textcomp}
\usepackage{xcolor}
\usepackage{hyperref}
\usepackage{booktabs}
\usepackage{listings}
\definecolor{codegreen}{rgb}{0,0.5,0}
\definecolor{codegray}{rgb}{0.4,0.4,0.4}
\definecolor{codepurple}{rgb}{0.5,0,0.5}
\lstdefinestyle{mystyle}{
    backgroundcolor=\color{white},   
    commentstyle=\color{codegreen},
    keywordstyle=\color{blue}\bfseries,
    numberstyle=\tiny\color{codegray},
    stringstyle=\color{codepurple},
    basicstyle=\ttfamily\tiny,
    breakatwhitespace=false,         
    breaklines=true,                 
    captionpos=b,                    
    keepspaces=true,                 
    numbers=left,                    
    numbersep=5pt,                  
    showspaces=false,                
    showstringspaces=false,
    showtabs=false,                  
    tabsize=2,
    frame=none
}
\lstset{style=mystyle}
\usepackage{float}
\usepackage{stfloats}
 


\renewcommand{\thesubsection}{\Alph{subsection}}
\newcommand{\tblsubhead}[1]{{\fontsize{7.5}{9}\selectfont #1}}

\title{Diseño e Implementación de un Brazo Robótico de 3 GDL para Corte Láser No-Coplanar}

\author{
\IEEEauthorblockN{\fontsize{9}{11}\selectfont Ivan Yerry Hinojosa Zegarra}
\IEEEauthorblockA{\fontsize{9}{11}\selectfont \textit{EAP Ingeniería Mecatrónica} \\
\textit{Universidad Continental}\\
Huancayo, Perú \\
72241950@continental.edu.pe}
\vspace{0.4cm} \\
\IEEEauthorblockN{\fontsize{9}{11}\selectfont Carlos Benedic Sosa Carhuamaca}
\IEEEauthorblockA{\fontsize{9}{11}\selectfont \textit{EAP Ingeniería Mecatrónica} \\
\textit{Universidad Continental}\\
Huancayo, Perú \\
60847160@continental.edu.pe}
\and
\IEEEauthorblockN{\fontsize{9}{11}\selectfont Brayan Omar Duran Arroyo}
\IEEEauthorblockA{\fontsize{9}{11}\selectfont \textit{EAP Ingeniería Mecatrónica} \\
\textit{Universidad Continental}\\
Huancayo, Perú \\
72542899@continental.edu.pe}
\and
\IEEEauthorblockN{\fontsize{9}{11}\selectfont Jorge Andre Tapia Patiño}
\IEEEauthorblockA{\fontsize{9}{11}\selectfont \textit{EAP Ingeniería Mecatrónica} \\
\textit{Universidad Continental}\\
Huancayo, Perú \\
60761518@continental.edu.pe}
\vspace{0.4cm} \\
\IEEEauthorblockN{\fontsize{9}{11}\selectfont Dick Jairo Guzman Zelada}
\IEEEauthorblockA{\fontsize{9}{11}\selectfont \textit{EAP Ingeniería Mecatrónica} \\
\textit{Universidad Continental}\\
Huancayo, Perú \\
60195546@continental.edu.pe}
}

\begin{document}

\maketitle

\begin{abstract}
    El presente artículo detalla el diseño, simulación y construcción física de un brazo manipulador antropomórfico de 3 grados de libertad (GDL) orientado al corte y perfilado láser no-coplanar. El objetivo principal ha sido desarrollar un sistema capaz de seguir trayectorias tridimensionales sobre superficies cilíndricas e irregulares mediante un emisor láser en su efector final, controlado por un microcontrolador ESP32. Se han implementado los modelos de cinemática directa, inversa y Jacobiana para garantizar el posicionamiento espacial en los tres ejes ($X, Y, Z$) y conservar la distancia focal del láser. La estructura física se fabricó en impresión 3D a partir de mallas STL, con un firmware desarrollado en C++. La validación cinemática se realizó mediante simulación virtual en ROS 2 y Gazebo de forma puramente offline.
\end{abstract}

\begin{IEEEkeywords}
Cinemática, brazo robótico, ESP32, C++, ROS 2, simulación virtual.
\end{IEEEkeywords}

\section{Introducción}
En la industria moderna, los brazos robóticos son herramientas fundamentales para tareas de precisión como la soldadura y el corte láser. Comprender la cinemática detrás de estos mecanismos es el primer paso para su control y programación \cite{b1}. Este proyecto busca aplicar los conceptos teóricos aprendidos en el curso de Fundamentos de Robótica ---tales como matrices de transformación homogénea, representación de orientación y cinemática diferencial--- en un caso de uso práctico. 

Al estructurar el modelo con 3 grados de libertad (GDL) (rotación en la base, elevación en el hombro y elevación en el codo), se simplificó el análisis geométrico sin perder la capacidad de posicionar el efector final en un volumen de trabajo tridimensional. Además, el proyecto utiliza un emisor láser de baja potencia (Clase 1/2) como efector final. A diferencia de procesos de contacto como la soldadura por arco, el láser permite una validación cinemática de ``no-contacto''. Esto simplifica el control al eliminar la necesidad de compensar fuerzas de fricción o colisiones mecánicas directas contra la pieza de trabajo, permitiendo un seguimiento de trayectoria preciso y eficiente. Este enfoque es crucial para el corte y grabado sobre superficies irregulares y curvadas en tres dimensiones, donde el brazo debe seguir contornos tridimensionales curvos sobre una superficie no plana manteniendo constante la distancia focal del haz láser, una operación inviable para máquinas cartesianas 2D tradicionales. El desarrollo del software, simulación, firmware embebido e interfaz web se encuentra disponible de forma pública en el repositorio de GitHub del proyecto \cite{b2}.

\section{Objetivos}
\subsection{Objetivo General}
Desarrollar, simular y construir un brazo robótico articulado de 3 GDL orientado a procesos de corte láser no-coplanar sobre superficies curvas e irregulares, validando los modelos cinemáticos en simulación virtual (ROS 2/Gazebo) y, de forma independiente, en el prototipo físico controlado por un ESP32.

\subsection{Objetivos Específicos}
\begin{itemize}
    \item Definir los parámetros de Denavit-Hartenberg y resolver analíticamente la cinemática directa, inversa y Jacobiana para garantizar el seguimiento espacial tridimensional ($X, Y, Z$) y la evasión de singularidades en superficies no planas.
    \item Modelar la estructura física en URDF/Xacro y simular el comportamiento del sistema utilizando el sistema operativo de robots (ROS 2, por sus siglas en inglés) y Gazebo de manera exclusivamente virtual, sin conexión física con el prototipo.
    \item Fabricar los eslabones y la base del robot utilizando tecnologías de impresión 3D a partir de archivos STL.
    \item Implementar el control de bajo nivel del hardware mediante un microcontrolador ESP32, utilizando C++ para accionar las articulaciones físicas del prototipo.
\end{itemize}

\section{Modelado Cinemático}
Se utiliza la convención de Denavit-Hartenberg (D-H) clásica para describir la geometría y cinemática del brazo manipulador antropomórfico de 3 GDL. 

\subsection{Sistemas de Referencia y Parámetros D-H}
Se definen cuatro sistemas de referencia locales $\{0\}$, $\{1\}$, $\{2\}$ y $\{3\}$ asociados a cada articulación móvil y al efector final de la estructura, coincidiendo con los marcos de simulación.
\begin{itemize}
    \item $d_1 = 60.000\text{ mm}$: Altura del eslabón base fijo.
    \item $d_2 = 36.173\text{ mm}$: Desplazamiento vertical del hombro.
    \item $a_2 = 14.915\text{ mm}$: Desplazamiento radial del hombro (distancia horizontal entre el eje de rotación de la base y el eje del hombro).
    \item $a_3 = 146.190\text{ mm}$: Longitud del primer eslabón móvil (distancia entre hombro y codo).
    \item $a_4 = 160.823\text{ mm}$: Longitud plana del segundo eslabón móvil (distancia de codo a brida/láser).
    \item $d_4 = -4.000\text{ mm}$: Desplazamiento transversal (fuera de plano) del efector final.
    \item $\theta_1$: Ángulo de rotación de la base (alrededor del eje vertical $Z_0$).
    \item $\theta_2$: Ángulo del hombro.
    \item $\theta_3$: Ángulo del codo.
\end{itemize}

La convención clásica de D-H establece que la matriz de transformación homogénea entre eslabones consecutivos se calcula mediante:
\begin{equation}
T_i^{i-1} = \begin{bmatrix} \cos\theta_i & -\sin\theta_i \cos\alpha_i & \sin\theta_i \sin\alpha_i & a_i \cos\theta_i \\ \sin\theta_i & \cos\theta_i \cos\alpha_i & -\cos\theta_i \sin\alpha_i & a_i \sin\theta_i \\ 0 & \sin\alpha_i & \cos\alpha_i & d_i \\ 0 & 0 & 0 & 1 \end{bmatrix}
\label{eq:dh_matrix}
\end{equation}

A partir de los sistemas de referencia y dimensiones físicas asignados, se obtiene la Tabla de Parámetros D-H (Tabla~\ref{tab:dh_parametros}). La Fig.~\ref{fig:dh_diagram} ilustra la disposición geométrica de estos parámetros sobre el modelo físico del brazo en su configuración de reposo ($\theta_1 = \theta_2 = \theta_3 = 0$).

\begin{table}[H]
\caption{Parámetros de Denavit-Hartenberg Reales del Prototipo}
\begin{center}
\footnotesize
\setlength{\tabcolsep}{4.5pt}
\begin{tabular}{c c c c c}
\toprule
\textbf{Eslabón ($i$)} & $\theta_i$ & $d_i$ \tblsubhead{(mm)} & $a_i$ \tblsubhead{(mm)} & $\alpha_i$ \\
\midrule
1 (0 $\rightarrow$ 1) & $\theta_1$ & $d_1 + d_2 = 96.173$ & $a_2 = 14.915$ & $90^\circ$ \\
2 (1 $\rightarrow$ 2) & $\theta_2$ & 0 & $a_3 = 146.190$ & 0 \\
3 (2 $\rightarrow$ 3) & $\theta_3$ & $d_4 = -4.000$ & $a_4 = 160.823$ & 0 \\
\bottomrule
\end{tabular}
\label{tab:dh_parametros}
\end{center}
\end{table}

\begin{figure}[!t]
\centering
\includegraphics[width=\columnwidth]{build/diagrama_p1.pdf}
\caption{Disposición geométrica del brazo robótico láser en configuración de reposo ($\theta_1\!=\!\theta_2\!=\!\theta_3\!=\!0$). Se indican los sistemas de referencia D-H $\{0\}$--$\{3\}$, las longitudes de eslabón $a_3$, $a_4$, los desplazamientos $d_1$, $d_2$ y el desfase transversal $d_4$.}
\label{fig:dh_diagram}
\end{figure}

\subsection{Cinemática Directa}
Sustituyendo los parámetros de la Tabla~\ref{tab:dh_parametros} en \eqref{eq:dh_matrix}, se obtienen las matrices de transformación homogénea individuales para cada articulación y eslabón:

\begin{equation}
T_1^0 = \begin{bmatrix} \cos\theta_1 & 0 & \sin\theta_1 & a_2\cos\theta_1 \\ \sin\theta_1 & 0 & -\cos\theta_1 & a_2\sin\theta_1 \\ 0 & 1 & 0 & d_1 + d_2 \\ 0 & 0 & 0 & 1 \end{bmatrix}
\end{equation}

\begin{equation}
T_2^1 = \begin{bmatrix} \cos\theta_2 & -\sin\theta_2 & 0 & a_3\cos\theta_2 \\ \sin\theta_2 & \cos\theta_2 & 0 & a_3\sin\theta_2 \\ 0 & 0 & 1 & 0 \\ 0 & 0 & 0 & 1 \end{bmatrix}
\end{equation}

\begin{equation}
T_3^2 = \begin{bmatrix} \cos\theta_3 & -\sin\theta_3 & 0 & a_4\cos\theta_3 \\ \sin\theta_3 & \cos\theta_3 & 0 & a_4\sin\theta_3 \\ 0 & 0 & 1 & d_4 \\ 0 & 0 & 0 & 1 \end{bmatrix}
\end{equation}

Multiplicando consecutivamente estas matrices se obtiene la matriz de transformación homogénea total del efector final respecto a la base del robot $T_3^0 = T_1^0 T_2^1 T_3^2$. Definimos esta matriz en función de las coordenadas del efector final $[x_E, y_E, z_E]^T$ dadas en \eqref{eq:fk_x}-\eqref{eq:fk_z}:
\begin{equation}
T_3^0 = \begin{bmatrix} C_1 C_{23} & -C_1 S_{23} & S_1 & x_E \\ S_1 C_{23} & -S_1 S_{23} & -C_1 & y_E \\ S_{23} & C_{23} & 0 & z_E \\ 0 & 0 & 0 & 1 \end{bmatrix}
\label{eq:matrix_total}
\end{equation}
donde se utiliza la notación compacta $C_1 = \cos\theta_1$, $S_1 = \sin\theta_1$, $C_{23} = \cos(\theta_2+\theta_3)$ y $S_{23} = \sin(\theta_2+\theta_3)$.

El vector de posición del efector final $P_E = [x_E, y_E, z_E]^T$ está dado por:
\begin{equation}
x_E = \cos\theta_1 \left( a_2 + a_3 \cos\theta_2 + a_4 \cos(\theta_2 + \theta_3) \right) + d_4 \sin\theta_1
\label{eq:fk_x}
\end{equation}
\begin{equation}
y_E = \sin\theta_1 \left( a_2 + a_3 \cos\theta_2 + a_4 \cos(\theta_2 + \theta_3) \right) - d_4 \cos\theta_1
\label{eq:fk_y}
\end{equation}
\begin{equation}
z_E = d_1 + d_2 + a_3 \sin\theta_2 + a_4 \sin(\theta_2 + \theta_3)
\label{eq:fk_z}
\end{equation}

Para verificar la consistencia del modelo matemático desarrollado, se evalúan las ecuaciones de cinemática directa bajo la configuración de reposo del prototipo (con todas las articulaciones en sus posiciones centrales $\theta_1 = 0$, $\theta_2 = 0$, $\theta_3 = 0$). En este estado, al sustituir los valores y las constantes físicas en \eqref{eq:fk_x}-\eqref{eq:fk_z}, se obtiene:
\begin{align}
x_E &= a_2 + a_3 \cos(0) + a_4 \cos(0) \nonumber \\
    &= 14.915 + 146.190 + 160.823 = 321.928\text{ mm} \\
y_E &= -d_4 = 4.000\text{ mm} \\
z_E &= d_1 + d_2 + a_3 \sin(0) + a_4 \sin(0) \nonumber \\
    &= 60.000 + 36.173 + 0 + 0 = 96.173\text{ mm}
\end{align}
Estos resultados analíticos coinciden exactamente con las coordenadas de la brida (efector final) obtenidas directamente del árbol de transformadas (TF) de la simulación virtual de ROS 2 ($X = 321.928\text{ mm}$, $Y = 4.000\text{ mm}$, $Z = 96.173\text{ mm}$). Esto valida la precisión del modelo cinemático de 4 sistemas propuesto para el guiado de trayectorias.

\subsection{Cinemática Inversa y Análisis de Múltiples Soluciones}
La resolución de la cinemática inversa consiste en determinar los ángulos físicos de los servomotores $(\theta_1, \theta_2, \theta_3)$ a partir de una posición cartesiana deseada $P_E = [x_E, y_E, z_E]^T$. Para ello, desacoplamos el efecto de la rotación de la base y del desfase transversal $d_4$:

1) \textit{Cálculo de $\theta_1$ (Rotación de la Base):} 
Proyectando el efector sobre el plano horizontal y compensando el desfase transversal $d_4$, definimos el radio proyectado plano $R_p$:
\begin{equation}
R_p = \sqrt{x_E^2 + y_E^2 - d_4^2}
\end{equation}
A partir del cual se calcula el ángulo de la base mediante:
\begin{equation}
\theta_1 = \operatorname{atan2}(y_E, x_E) - \operatorname{atan2}(-d_4, R_p)
\label{eq:ik_theta1}
\end{equation}

2) \textit{Cálculo de $\theta_3$ (Elevación del Codo):}
Se define la proyección en el plano del brazo. El radio proyectado corregido $r$ y la altura corregida respecto al hombro $z'$ se formulan como:
\begin{equation}
r = R_p - a_2
\end{equation}
\begin{equation}
z' = z_E - (d_1 + d_2)
\end{equation}
Aplicando la ley de los cosenos en el triángulo del plano del brazo para resolver el ángulo $\theta_3$:
\begin{equation}
r^2 + z'^2 = a_3^2 + a_4^2 + 2 a_3 a_4 \cos\theta_3
\end{equation}
Despejando el término del coseno se obtiene:
\begin{equation}
\cos\theta_3 = \frac{r^2 + z'^2 - a_3^2 - a_4^2}{2 a_3 a_4} = \Psi
\label{eq:ik_cos3}
\end{equation}
Para la existencia de una solución física real, se debe cumplir $-1 \le \Psi \le 1$. La ecuación admite dos soluciones:
\begin{equation}
\theta_3 = \operatorname{atan2}\left( \pm\sqrt{1 - \Psi^2}, \Psi \right)
\label{eq:ik_theta3}
\end{equation}
donde el signo positivo ($+$) representa la configuración de \textbf{codo abajo} y el signo negativo ($-$) representa la de \textbf{codo arriba} (flexión natural hacia abajo en el modelo calibrado). Seleccionamos la configuración de \textbf{codo arriba} para evitar colisiones mecánicas del brazo con la mesa de trabajo.

3) \textit{Cálculo de $\theta_2$ (Elevación del Hombro):}
Una vez obtenido $\theta_3$, se resuelve para $\theta_2$ definiendo las constantes auxiliares $k_1 = a_3 + a_4 \cos\theta_3$ y $k_2 = a_4 \sin\theta_3$:
\begin{equation}
\theta_2 = \operatorname{atan2}\left( k_1 z' - k_2 r, k_1 r + k_2 z' \right)
\label{eq:ik_theta2}
\end{equation}

\subsection{Jacobiano y Análisis de Singularidades}
El Jacobiano lineal del manipulador $J_v \in \mathbb{R}^{3 \times 3}$ relaciona las velocidades articulares $\dot{q} = [\dot{\theta}_1, \dot{\theta}_2, \dot{\theta}_3]^T$ con la velocidad lineal del efector final. Definimos el radio horizontal proyectado $R_p = a_2 + a_3 \cos\theta_2 + a_4 \cos(\theta_2 + \theta_3)$. Derivando parcialmente las ecuaciones de la cinemática directa \eqref{eq:fk_x}-\eqref{eq:fk_z}:
\begin{equation}
J_v = \begin{bmatrix} -S_1 R_p + d_4 C_1 & -C_1(a_3 S_2 + a_4 S_{23}) & -C_1 a_4 S_{23} \\ C_1 R_p + d_4 S_1 & -S_1(a_3 S_2 + a_4 S_{23}) & -S_1 a_4 S_{23} \\ 0 & a_3 C_2 + a_4 C_{23} & a_4 C_{23} \end{bmatrix}
\label{eq:jacobian}
\end{equation}
donde $S_i$, $C_i$, $S_{23}$ y $C_{23}$ representan la notación simplificada de los senos y cosenos directos.

Calculando analíticamente el determinante de $J_v$ (despreciando el término menor de desfase $d_4$):
\begin{equation}
\det(J_v) = -a_3 a_4 \left( a_2 + a_3 \cos\theta_2 + a_4 \cos(\theta_2 + \theta_3) \right) \sin\theta_3
\label{eq:det_jacobian}
\end{equation}

Se identifican las dos configuraciones singulares críticas:
\begin{itemize}
    \item \textbf{Singularidad del codo ($\sin\theta_3 = 0$):} Ocurre cuando $\theta_3 = 0$ o $\theta_3 = -\pi$ (brazo completamente estirado o retraído).
    \item \textbf{Singularidad de alineación con el eje de la base ($R_p = 0$):} Ocurre cuando el efector final coincide con el eje de rotación de la base ($Z_0$).
\end{itemize}




\section{Metodología y Alcance}
La ejecución del proyecto se estructuró en dos fases metodológicas consecutivas:

\subsection{Fase 1: Modelado Mecánico y Simulación Virtual}
Se diseñaron las piezas tridimensionales del manipulador en CAD. Los modelos tridimensionales de los eslabones se exportaron en formato STL para el proceso de manufactura, y se compilaron en un archivo descriptor URDF (con formato de mallas colisionables e inerciales). En el entorno de ROS 2 y Gazebo, se programó un nodo de simulación independiente para validar computacionalmente los modelos de cinemática directa e inversa, garantizando que el seguimiento de las trayectorias cartesianas de corte ocurra libre de colisiones mecánicas directas o singularidades. Esta simulación en ROS 2 y Gazebo se mantiene de forma puramente virtual, sin establecer conexiones directas de comunicación o control hacia el prototipo físico.

\subsection{Fase 2: Fabricación e Integración de Hardware}
La estructura mecánica se fabricó mediante impresión 3D FDM de los archivos STL utilizando plástico PETG con un patrón de relleno de alta densidad (40\% o superior) para absorber esfuerzos de flexión. Se integraron los actuadores en sus correspondientes articulaciones físicas: un servomotor MG996R en la base (giro izquierdo/derecha), un servomotor MG996R en el hombro (para soportar el torque gravitacional) y un servomotor MG90S en el codo.


\section{Implementación de Software y Algoritmos de Control}
El ecosistema de software del manipulador se divide en tres capas fundamentales que operan coordinadamente de forma puramente virtual para el entorno de simulación (la interfaz web interactiva de usuario, el generador síncrono de trayectorias en ROS 2 y el puente de comunicación de red):

\subsection{Arquitectura de Comunicación e Intercambio de Datos}
El ecosistema de software opera en un esquema coordinado donde la interfaz de usuario interactiva y el resolvedor dinámico de la simulación física se comunican en tiempo real a través de un puente HTTP y tópicos de ROS 2. El flujo de datos e integración de los componentes clave se describe a continuación:
\begin{itemize}
    \item \textbf{Capa de Interfaz Gráfica (Frontend):} Implementada en HTML/JS (\texttt{app.js}), captura la interacción del operador mediante barras de desplazamiento para los ángulos articulares o entradas cartesianas tridimensionales de posición ($X, Y, Z$) calculadas localmente mediante cinemática inversa. Cuando ocurre un cambio, la función \texttt{sendJointsToROS()} serializa los ángulos de las articulaciones $J_1, J_2, J_3$ (en grados sexagesimales) dentro de un objeto JSON y realiza una petición HTTP POST throttleada (con un intervalo mínimo de $50\text{ ms}$) hacia el endpoint local de la API \texttt{/api/move}.
    \item \textbf{Capa del Puente de Red (Web Bridge):} El script en Python \texttt{web\_bridge.py} inicializa de forma concurrente un servidor HTTP en el puerto $8080$ y un nodo de ROS 2. Al recibir la petición POST en la ruta \texttt{/api/move}, el manejador de peticiones del servidor decodifica el objeto JSON, recupera los valores angulares y ejecuta la función \texttt{publish\_joints()} del nodo. Esta función convierte de forma instantánea las variables articulares de grados sexagesimales a radianes y las publica en el tópico \texttt{/arm\_controller/commands} de ROS 2.
    \item \textbf{Capa de Simulación Virtual (ROS 2 / Gazebo):} El resolvedor físico y dinámico de Gazebo está suscrito de forma nativa a \texttt{/arm\_controller/commands}. Al recibir los radianes de consigna, los controladores de posición articulares virtuales aplican torques proporcionales para llevar los eslabones simulados del robot a la pose deseada.
    \item \textbf{Generación de Trayectorias Autónomas:} De manera paralela e independiente de la interfaz gráfica, el script \texttt{trajectory\_generator.py} publica directamente en el tópico \texttt{/arm\_controller/commands} a una frecuencia de $50\text{ Hz}$. Este nodo interpola entre una secuencia programada de keyframes utilizando un perfil de velocidad cosenoidal (S-curve), permitiendo realizar validaciones dinámicas automáticas dentro del simulador físico Gazebo de forma desacoplada del frontend web.
\end{itemize}

\subsection{Interfaz Web de Control e Interacción Dinámica (app.js)}
La interfaz de usuario está desarrollada en JavaScript de cliente puro y HTML5. Permite realizar simulación en línea y resolvedor de cinemática inversa interactiva. El algoritmo detallado de la cinemática inversa se presenta en el Código~\ref{lst:ik_js} en el Apéndice~\ref{ap:app_js}.
Los componentes de la capa de interfaz web realizan las siguientes tareas:
\begin{itemize}
    \item \textbf{Lienzo de Dibujo Interactivo (Canvas):} La rutina \texttt{drawRobot()} proyecta en tiempo real el comportamiento dinámico bidimensional del robot sobre un lienzo Canvas HTML5, mostrando vistas superior ($X$-$Y$) y lateral ($X$-$Z$) para dar retroalimentación visual al operador sobre la trayectoria de corte.
    \item \textbf{Interpolación de Suavizado de Movimiento:} Para transicionar sin saltos discretos hacia las coordenadas indicadas, la función \texttt{animateToJoints()} aplica un algoritmo de facilidad cúbica (\textit{Cubic Ease-Out}) a $60\text{ cuadros por segundo}$, emulando el amortiguamiento físico de los servomotores reales.
\end{itemize}

\subsection{Generador de Trayectorias de Simulación (trajectory\_generator.py)}
El nodo generador de trayectorias en Python modela el perfil cinemático de interpolación de simulación en ROS 2 y Gazebo. El lazo periódico de control temporal de este script se detalla en el Código~\ref{lst:traj_py} en el Apéndice~\ref{ap:traj_py}.
Este nodo realiza la generación y control síncrono mediante:
\begin{itemize}
    \item \textbf{Interpolación Cosenoidal (S-curve):} La ecuación cosenoidal suaviza el incremento de velocidad al inicio y finalización del movimiento de cada eslabón, eliminando las aceleraciones abruptas que provocan oscilaciones inerciales o colisiones mecánicas virtuales en el resolvedor de física de Gazebo.
    \item \textbf{Bucle Síncrono a $50\text{ Hz}$:} Garantiza una comunicación continua hacia los controladores articulares virtuales a intervalos exactos de $20\text{ ms}$.
\end{itemize}

\subsection{Puente de Comunicación ROS 2 - Interfaz Web (web\_bridge.py)}
La comunicación intermedia entre la simulación y la interfaz de usuario web interactiva se gestiona mediante un puente de traducción bidireccional desarrollado en Python. La conversión de unidades y publicación de comandos se realiza a través del método detallado en el Código~\ref{lst:bridge_py} en el Apéndice~\ref{ap:bridge_py}.
Este puente opera sobre las siguientes bases técnicas:
\begin{itemize}
    \item \textbf{Servidor HTTP Multihilo:} Inicia un servidor web HTTP concurrente en el puerto 8080. La interfaz de usuario realiza peticiones de tipo POST con formato JSON hacia la ruta `/api/move`.
    \item \textbf{Traducción de Coordenadas Articulares:} Convierte las entradas angulares desde el formato en grados sexagesimales usado por la interfaz web al formato nativo en radianes esperado por los controladores virtuales de ROS 2, garantizando la sintonía espacial entre la simulación interactiva y el gemelo virtual.
\end{itemize}

\subsection{Procesamiento de Imagen y Generación de Trayectoria 3D (web\_bridge.py)}
Para permitir que el brazo robótico realice trayectorias de corte basadas en dibujos o patrones de usuario, se implementó una canalización (pipeline) de visión artificial en el puente de comunicación:
\begin{itemize}
    \item \textbf{Esqueletización de Imagen (Zhang-Suen):} Reduce los trazos de dibujo binarios de la imagen subida a líneas continuas de un solo píxel de espesor utilizando una implementación vectorizada en NumPy del algoritmo de adelgazamiento de Zhang-Suen, detallada en el Código~\ref{lst:zhang_py} en el Apéndice~\ref{ap:zhang_py}. Esto optimiza drásticamente el tiempo de procesamiento en comparación con bucles tradicionales.
    \item \textbf{Proyección a Coordenadas Cartesianas 3D:} Los píxeles esqueletizados en 2D se proyectan al plano tridimensional utilizando el modelo geométrico de la superficie de trabajo cilíndrica. Cada píxel se asocia a un punto tridimensional $(x, y, z)$ en metros y almacena la instrucción de encendido del emisor láser.
\end{itemize}

\subsection{Orquestación y Lanzamiento de Simulación (gazebo.launch.py)}
La configuración y el inicio coordinado de los nodos de simulación física en ROS 2 se gestionan mediante un archivo de lanzamiento desarrollado en Python. El flujo secuencial y la configuración del entorno se detallan en el Código~\ref{lst:launch_py} en el Apéndice~\ref{ap:launch_py}.
Este orquestador realiza las siguientes operaciones:
\begin{itemize}
    \item \textbf{Carga y Procesamiento Dinámico de URDF/Xacro:} El script utiliza la utilidad \texttt{xacro.process\_file()} para interpretar el modelo XML parametrizado. Posteriormente, aplica expresiones regulares para remover los comentarios XML redundantes con el fin de prevenir errores en el parser de ROS 2 y evitar inconsistencias en el plugin de Gazebo.
    \item \textbf{Gestión de Rutas de Modelos de Gazebo:} Actualiza dinámicamente la variable de entorno \texttt{GAZEBO\_MODEL\_PATH} para incluir el directorio local de mallas STL del brazo, garantizando la carga correcta de las geometrías de visualización en Gazebo Classic.
    \item \textbf{Secuenciación de Nodos mediante Manejo de Eventos:} Para evitar condiciones de carrera donde los controladores articulares intentan enlazarse antes de que el robot exista físicamente en la física del simulador, el script utiliza manejadores de eventos \texttt{RegisterEventHandler} y \texttt{OnProcessExit}. Esto fuerza que el controlador de posición de las articulaciones (\texttt{arm\_controller}) se lance únicamente después de que el robot haya sido instanciado por el nodo \texttt{spawn\_entity}.
\end{itemize}

\subsection{Modelo Descriptor de la Estructura y Actuadores (arm.urdf.xacro)}
El diseño estructural, los límites de movimiento y la configuración de hardware de control del gemelo virtual se definen jerárquicamente en un archivo descriptor Xacro, presentado en el Código~\ref{lst:urdf_xacro} en el Apéndice~\ref{ap:urdf_xacro}.
Este descriptor se estructura a partir de los siguientes elementos clave:
\begin{itemize}
    \item \textbf{Geometría de Eslabones y Mallas STL:} Vincula directamente los archivos STL exportados de CAD (\texttt{base.stl}, \texttt{link1.stl}, etc.) escalados a metros (factor $0.001$), definiendo la apariencia visual del robot en el simulador.
    \item \textbf{Límites y Dinámica de Articulaciones:} Modela los límites físicos angulares en radianes equivalentes a los límites de los servomotores físicos del prototipo. Adicionalmente, incluye parámetros dinámicos de fricción y amortiguamiento (\texttt{friction} y \texttt{damping}) para modelar el rozamiento de las articulaciones en Gazebo.
    \item \textbf{Integración con ros2\_control:} Declara el elemento \texttt{<ros2\_control>} que mapea las interfaces de comando y de estado (posiciones articulares y velocidades) del resolvedor de física de Gazebo Classic (\texttt{gazebo\_ros2\_control/GazeboSystem}), permitiendo la inyección directa de señales desde el puente de comunicación HTTP. Esta declaración se complementa en la arquitectura de ROS 2 con el archivo de configuración \texttt{controllers.yaml}, el cual especifica la tasa de actualización de los bucles de control a $100\text{ Hz}$ y define un controlador síncrono del tipo \texttt{JointGroupPositionController} asociado a las tres articulaciones activas del robot.
\end{itemize}

\subsection{Automatización de Configuración y Compilación (setup.sh)}
Para simplificar la inicialización del entorno y garantizar la reproducibilidad de las pruebas, se programó un script de automatización en Bash. El funcionamiento y la lógica de este script se detallan en el Código~\ref{lst:setup_sh} en el Apéndice~\ref{ap:setup_sh}.
Este script realiza las siguientes operaciones secuenciales:
\begin{itemize}
    \item \textbf{Verificación del Entorno Aislado:} Comprueba la existencia del contenedor Distrobox (\texttt{ros2-humble}) y, en caso de no encontrarse activo, lo crea automáticamente instalando una base de Ubuntu 22.04.
    \item \textbf{Compilación Síncrona del Workspace:} Accede al contenedor, carga las variables de entorno de ROS 2 Humble y compila exclusivamente el paquete \texttt{arm\_simulation} mediante el comando \texttt{colcon build}.
    \item \textbf{Generación de Scripts Lanzadores:} Escribe y otorga permisos de ejecución a los tres scripts rápidos (\texttt{launch\_gazebo.sh}, \texttt{run\_trajectory.sh} y \texttt{start\_ui.sh}) en la raíz del proyecto para facilitar el uso diario.
\end{itemize}

\subsection{Scripts Lanzadores del Entorno y Ejecutables}
Para agilizar el flujo de trabajo del operador sin requerir la memorización de comandos de contenedor o de ROS 2, se desarrollaron tres scripts lanzadores. Su código fuente se presenta en el Código~\ref{lst:runners_sh} en el Apéndice~\ref{ap:runners_sh}.
Cada script cumple una función específica dentro de la arquitectura distribuida:
\begin{itemize}
    \item \textbf{Lanzamiento de Simulación (launch\_gazebo.sh):} Limpia la variable de entorno \texttt{PATH} para remover rutas de Conda (previniendo conflictos de intérprete de Python con las dependencias de ROS 2), carga los entornos locales y ejecuta el launch de Gazebo.
    \item \textbf{Generador de Trayectorias (run\_trajectory.sh):} Realiza la misma sanitización del entorno y ejecuta directamente el nodo controlador de interpolación cosenoidal.
    \item \textbf{Inicio de la Interfaz Gráfica (start\_ui.sh):} Lanza de manera asíncrona el navegador predeterminado apuntando a \texttt{http://localhost:8080} y arranca concurrentemente el puente HTTP/ROS 2 (\texttt{web\_bridge.py}) en segundo plano.
\end{itemize}

\subsection{Firmware Embebido del Microcontrolador Físico (main.cpp)}
Para el control directo y la coordinación de los actuadores del prototipo físico, se desarrolló un firmware embebido en C++ ejecutado sobre el microcontrolador ESP32. Su implementación completa se detalla en el Código~\ref{lst:main_cpp} en el Apéndice~\ref{ap:main_cpp}. Las principales funcionalidades de control y seguridad implementadas se describen a continuación:
\begin{itemize}
    \item \textbf{Configuración de Periféricos e Inicialización de Actuadores:} El firmware utiliza la biblioteca \texttt{ESP32Servo} reservando tres temporizadores de hardware a una frecuencia de control de $50\text{ Hz}$ para generar las señales de modulación por ancho de pulso (PWM) destinadas a los servomotores de la base (\texttt{GPIO 27}), hombro (\texttt{GPIO 14}) y codo (\texttt{GPIO 21}). Adicionalmente, el pin \texttt{GPIO 22} se define como una salida digital para la conmutación del emisor láser.
    \item \textbf{Resolvedor de Cinemática Inversa en Tiempo Real:} El microcontrolador ejecuta localmente el modelo matemático inverso analítico codo-arriba descrito en las ecuaciones \eqref{eq:ik_theta1}-\eqref{eq:ik_theta2}. A partir de una coordenada cartesiana de destino $(x, y, z)$, calcula dinámicamente los ángulos en radianes $(\theta_1, \theta_2, \theta_3)$ y los traduce a microsegundos de ancho de pulso mediante la función \texttt{angleToMicroseconds()}, restringiendo los valores entre los límites físicos seguros de $[500, 2500]\,\mu\text{s}$.
    \item \textbf{Algoritmo de Movimiento Gradual Simultáneo:} Al comandar saltos instantáneos de posición, los servomotores demandan picos elevados de corriente instantánea que activan la protección contra sobrecorriente (OCP) de la fuente de alimentación externa, provocando el apagado del sistema. Para resolver este problema físico, se implementó el método \texttt{moveToTarget()}. Este algoritmo desplaza de forma síncrona y paralela los tres actuadores en incrementos máximos de $15\,\mu\text{s}$ ($\approx 1.35^\circ$) espaciados por retardos temporales de $15\text{ ms}$. Esto limita la tasa de variación de la corriente eléctrica durante los arranques y detenciones de las articulaciones.
    \item \textbf{Decodificación de Tramas y Auto-escalado de Unidades:} El bucle de control analiza las tramas provenientes de la interfaz serial esperando el formato de texto plano \texttt{"X:val,Y:val,Z:val"}. Con el fin de flexibilizar la operación, el programa comprueba si las coordenadas ingresadas están expresadas en milímetros (valores absolutos mayores a $2.0$). Si es el caso, aplica automáticamente un factor de escala de $0.001$ para convertirlas a metros antes de realizar el cálculo cinemático.
    \item \textbf{Mecanismos de Seguridad Integrados:} Para prevenir quemaduras incidentales o incendios en la pieza de trabajo por pérdida de conexión con la interfaz de control, el código implementa un temporizador de watchdog por software. Si transcurren más de $1000\text{ ms}$ sin recibir comandos seriales válidos, el emisor láser se desactiva inmediatamente (\texttt{digitalWrite(LASER\_PIN, LOW)}). Asimismo, la rutina \texttt{triggerEmergency()} permite interrumpir de forma inmediata la señal PWM y apagar el láser ante la recepción de tramas manuales de emergencia como \texttt{ESTOP} o \texttt{EMERGENCY}.
\end{itemize}

\subsection{Scripts de Validación Cinemática y Alineación de Marcos}
Con el fin de verificar la corrección del URDF, el árbol de transformaciones (TF) y las ecuaciones cinemáticas antes de operar el hardware, se programaron diversos scripts de prueba en Python:
\begin{itemize}
    \item \textbf{Prueba Cinemática y Comparación DH (test\_kinematics.py):} Compara analíticamente dos formulaciones de Denavit-Hartenberg (una estándar y otra con desfases angulares compensados $\phi_2$ y $\phi_3$), validando la consistencia matemática de la cinemática directa. El código se detalla en el Código~\ref{lst:test_kin} en el Apéndice~\ref{ap:test_kin}.
    \item \textbf{Validación de Alineación TF en Gazebo (test\_tf\_alignment.py):} Se conecta como nodo de ROS 2 para leer las transformaciones del árbol de TFs en tiempo real de la simulación activa. Comprueba que en la pose cero de calibración del robot, todos los ejes longitudinales $X$ de las articulaciones sean coplanares y paralelos (Ángulo Pitch $\approx 0^\circ$ y desviación lateral $Y \approx 0$). Su código se presenta en el Código~\ref{lst:test_tf} en el Apéndice~\ref{ap:test_tf}.
    \item \textbf{Verificación Analítica de URDF (test\_urdf\_tf\_alignment.py):} Realiza un análisis algebraico de las transformaciones de rotación RPY declaradas en el modelo descriptor URDF/Xacro, comprobando si los ejes longitudinales $X$ resultan paralelos y horizontales en el plano de movimiento vertical. Su código se detalla en el Código~\ref{lst:test_urdf_tf} en el Apéndice~\ref{ap:test_urdf_tf}.
    \item \textbf{Lectura de Marcos de Referencia Activos (get\_tf\_distances.py):} Consulta el servicio de transformaciones de ROS 2 en tiempo real para imprimir las traslaciones y cuaterniones de rotación de cada eslabón con respecto a la base del robot. Su implementación se muestra en el Código~\ref{lst:get_tf_dist} en el Apéndice~\ref{ap:get_tf_dist}.
    \item \textbf{Parser de Geometría Xacro (get\_gazebo\_distances.py):} Analiza el árbol XML del archivo Xacro descriptor del brazo para extraer automáticamente los límites físicos, desfases de juntas y longitudes calculadas de los eslabones móviles. Su código se detalla en el Código~\ref{lst:get_gz_dist} en el Apéndice~\ref{ap:get_gz_dist}.
\end{itemize}

\section{Selección de Actuadores y Sensores}
Para garantizar el éxito de la implementación física y su congruencia con el entorno de simulación, se detalla la selección de componentes de hardware:


\subsection{Actuadores}
La arquitectura cinemática de 3 GDL resulta idónea para este caso de estudio; dado que la herramienta es un haz láser unidireccional colineal con el eje de aproximación, y el objetivo es el corte y grabado sobre superficies cilíndricas, el control del posicionamiento cartesiano en los tres ejes ($X, Y, Z$) es crítico para compensar la curvatura de la pieza y mantener la distancia focal constante del haz. La omisión del control de orientación redundante (pitch/roll) disminuye drásticamente el torque estático requerido en las articulaciones proximales (base y hombro), lo que mitiga las deflexiones estructurales y maximiza la rigidez y repetibilidad del sistema. El diseño se complementa ubicando estratégicamente actuadores de baja masa en los eslabones distales para minimizar el tensor de inercia del mecanismo.

\begin{itemize}
    \item \textbf{Actuador de rotación principal (Base):} Se implementó un servomotor de alto torque y engranajes metálicos (MG996R) en la base para accionar el giro del brazo a la izquierda y derecha. Este actuador soporta las cargas dinámicas y el momento de inercia de toda la estructura móvil.
    \item \textbf{Actuadores de elevación (Hombro y Codo):} Se utilizó un servomotor de alto torque MG996R para el hombro (articulación pegada a la base que hace elevar el brazo) para soportar el peso y los esfuerzos de flexión de los eslabones, y un servomotor de engranajes metálicos MG90S en la articulación del codo para controlar la elevación distal.
    \item \textbf{Efector final (Módulo Láser de estado sólido):} Consiste en un diodo láser de baja potencia con lógica transistor-transistor (TTL) de 5V, el cual es conmutado como un nivel lógico de entrada/salida de propósito general (GPIO) o vía PWM desde el ESP32. Esto permite una activación síncrona y de baja latencia acoplada al perfil cinemático de la trayectoria, emulando la dinámica de interpolación de un proceso industrial de corte láser CNC.
\end{itemize}

\subsection{Sensores}
\begin{itemize}
    \item \textbf{Potenciómetros internos:} Los servomotores ya incluyen sensores de posición (potenciómetros) internamente. Esto mantiene el diseño mecánico limpio ya que no se requiere acoplar encoders ópticos externos para leer los ángulos de las articulaciones, eliminando la necesidad de cableado para sensores externos de posición o corriente.
\end{itemize}

\subsection{Especificación de Componentes}
En la Tabla~\ref{tab:componentes} se detallan las especificaciones y funciones de los componentes de hardware seleccionados.

\begin{table}[H]
\caption{Especificaciones de los Componentes de Hardware}
\begin{center}
\footnotesize
\begin{tabular}{l l p{3.0cm}}
\toprule
\textbf{Componente} & \textbf{Modelo} & \textbf{Función en el Sistema} \\
\midrule
Controlador & ESP32 & Control de bajo nivel y PWM. \\
Actuador base & MG996R (1 ud.) & Rotación izquierda/derecha. \\
Actuador hombro & MG996R (1 ud.) & Elevación del hombro. \\
Actuador codo & MG90S (1 ud.) & Elevación del codo. \\
Efector final & Diodo Láser (5V) & Haz de corte de baja potencia. \\
Estructura & PETG & Piezas impresas en 3D. \\
Alimentación & Fuente 5V 5A & Energía para actuadores y láser. \\
\bottomrule
\end{tabular}
\label{tab:componentes}
\end{center}
\end{table}

\subsection{Presupuesto de Masa y Dimensiones}
En la Tabla~\ref{tab:masa_dimensiones} se detalla la masa y dimensiones físicas estimadas para los diferentes componentes del brazo manipulador.

\begin{table}[H]
\caption{Presupuesto de Masa y Dimensiones de los Componentes}
\begin{center}
\footnotesize
\begin{tabular}{l c c p{2.8cm}}
\toprule
\textbf{Componente} & \textbf{Masa} \tblsubhead{(g)} & \textbf{Cant.} & \textbf{Dimensiones} \tblsubhead{(mm)} \\
\midrule
ESP32 NodeMCU & 10 & 1 & 48.0 $\times$ 26.0 $\times$ 12.0 \\
Servo MG996R & 55 & 2 & 40.7 $\times$ 19.7 $\times$ 42.9 \\
Servo MG90S & 13.4 & 1 & 22.8 $\times$ 12.2 $\times$ 28.5 \\
Diodo Láser 5V & 10 & 1 & $\phi$ 12.0 $\times$ 35.0 \\
Estructura 3D (PETG) & 350 & 1 & Envolvente $\approx$ 250 $\times$ 150 \\
Fuente 5V 5A & 150 & 1 & 110.0 $\times$ 78.0 $\times$ 36.0 \\
Cables y tornillería & 50 & 1 & N/A \\
\midrule
\textbf{Masa Total} & \textbf{693 g} & \multicolumn{2}{l}{\tblsubhead{($\approx$ 0.69 kg)}} \\
\bottomrule
\end{tabular}
\label{tab:masa_dimensiones}
\end{center}
\end{table}

\subsection{Presupuesto y Consumo de Energía}
La Tabla~\ref{tab:energia} detalla el voltaje nominal y las corrientes de funcionamiento promedio y de pico para estimar la potencia del sistema.

\begin{table}[H]
\caption{Presupuesto y Consumo de Energía de los Componentes}
\begin{center}
\footnotesize
\begin{tabular}{l c c c}
\toprule
\textbf{Componente} & \textbf{Voltaje} \tblsubhead{(V)} & \textbf{Corr. Prom.} & \textbf{Corr. Pico} \\
\midrule
ESP32 & 5.0 & 120 mA & 240 mA \\
MG996R (2 ud.) & 5.0 & 1000 mA & 5.0 A \\
MG90S (1 ud.) & 5.0 & 180 mA & 0.8 A \\
Láser 5V & 5.0 & 30 mA & 30 mA \\
\midrule
\multicolumn{2}{l}{\textbf{Consumo Promedio}} & \multicolumn{2}{r}{\textbf{1.33 A} \tblsubhead{($\approx$ 6.7 W)}} \\
\multicolumn{2}{l}{\textbf{Consumo Pico}} & \multicolumn{2}{r}{\textbf{6.07 A} \tblsubhead{($\approx$ 30.4 W)}} \\
\bottomrule
\end{tabular}
\label{tab:energia}
\end{center}
\end{table}

\subsection{Estimación de Costos}
El presupuesto estimado para la adquisición de los componentes de hardware requeridos se presenta en la Tabla~\ref{tab:costos}.

\begin{table}[H]
\caption{Presupuesto y Estimación de Costos}
\begin{center}
\footnotesize
\begin{tabular}{l c r r}
\toprule
\textbf{Descripción} & \textbf{Cant.} & \textbf{Unit.} \tblsubhead{(S/.)} & \textbf{Total} \tblsubhead{(S/.)} \\
\midrule
ESP32 NodeMCU & 1 & 25.00 & 25.00 \\
Servomotor MG996R & 2 & 30.00 & 60.00 \\
Servomotor MG90S & 1 & 12.00 & 12.00 \\
Módulo Láser 5V TTL & 1 & 15.00 & 15.00 \\
Filamento PETG (1 kg) & 1 & 70.00 & 70.00 \\
Fuente de alimentación 5V 5A & 1 & 35.00 & 35.00 \\
Cables y tornillería & - & 15.00 & 15.00 \\
\midrule
\multicolumn{3}{l}{\textbf{Total Estimado}} & \textbf{S/. 232.00} \\
\bottomrule
\end{tabular}
\label{tab:costos}
\end{center}
\end{table}

\begin{figure}[!b]
\centering
\includegraphics[width=\columnwidth]{build/diagrama_p2.pdf}
\caption{Esquema de conexiones eléctricas del prototipo. Se detallan las líneas de alimentación desde la fuente de PC (+5\,V y GND) hacia los servomotores Servo~1 (MG996R), Servo~2 (MG996R) y Servo~3 (MG90S), el circuito de control del diodo láser con transistor NPN 2N2222 y resistencia de base de 220\,$\Omega$, y la conexión USB de alimentación del ESP32 NodeMCU de 38 pines.}
\label{fig:electrical_diagram}
\end{figure}

\subsection{Esquema de Conexiones Eléctricas}
La integración física del sistema requiere una arquitectura de alimentación y control claramente definida. El ESP32 es alimentado directamente por la laptop vía USB, mientras que los servomotores y el diodo láser reciben su tensión de operación de +5\,V desde la fuente de PC. El encendido del láser se gestiona mediante un transistor NPN 2N2222 con resistencia de base de 220\,$\Omega$, accionado por el pin GPIO~25 del ESP32. La Fig.~\ref{fig:electrical_diagram} muestra el esquema completo de conexiones.

\section{Aplicación Industrial}
A diferencia de los pantógrafos y cortadoras láser cartesianas convencionales, los cuales están limitados a trazar trayectorias coplanares bidimensionales ($2\text{D}$) en láminas planas, la configuración antropomórfica de 3 GDL permite posicionar la herramienta en un espacio tridimensional ($X, Y, Z$). Esta característica habilita la aplicación del robot en procesos industriales de corte y grabado sobre superficies irregulares y no coplanares, tales como el perfilado de cuerpos curvados y el calado de piezas cilíndricas.

En la industria de manufactura mecánica, el procesamiento de piezas no planas requiere realizar cortes y grabados con perfiles tridimensionales complejos (por ejemplo, intersecciones de geometrías cilíndricas o marcas en piezas curvas). Al colocar un haz láser en el efector final del robot, es posible realizar este proceso sin contacto mecánico directo. A partir del modelo de la superficie cilíndrica de una pieza de radio $R_{tubo}$ colocada de forma horizontal a lo largo del eje $Y_0$ a una distancia $x_{tubo}$ y a una altura de desfase $z_{tubo}$, la coordenada vertical de corte $z_E$ se calcula dinámicamente en función de la posición radial $x_E$:
\begin{equation}
z_E = z_{tubo} + \sqrt{R_{tubo}^2 - (x_E - x_{tubo})^2}
\label{eq:tubo}
\end{equation}
El resolvedor de cinemática inversa desarrollado en \eqref{eq:ik_theta1}-\eqref{eq:ik_theta2} calcula de manera instantánea las posiciones articulares $(\theta_1, \theta_2, \theta_3)$ para que el láser siga esta trayectoria no planar, manteniendo la distancia focal constante respecto a la pared de la pieza. Esto demuestra el valor industrial del manipulador propuesto frente a las soluciones cartesianas convencionales de bajo costo, ampliando el alcance de automatización hacia geometrías complejas no convencionales.

\appendices

\section{Código del Resolvedor de Cinemática Inversa en la Interfaz Web (app.js)}
\label{ap:app_js}
\begin{lstlisting}[language=C++,caption={\small Resolvedor interactivo de cinemática inversa en JS (app.js).},label={lst:ik_js}]
function solveIK(x, y, z) {
  const targetX = parseFloat(x);
  const targetY = parseFloat(y);
  const targetZ = parseFloat(z);

  if (isNaN(targetX) || isNaN(targetY) || isNaN(targetZ)) {
    return { error: 'VALORES_INVALIDOS', msg: 'Ingresa valores numericos validos.' };
  }

  let r = Math.sqrt(targetX * targetX + targetY * targetY);
  if (r < Math.abs(d4)) {
    r = Math.abs(d4);
  }
  const Rp = Math.sqrt(r * r - d4 * d4);
  const theta1 = (targetX === 0 && targetY === 0) ? -Math.PI / 2 : Math.atan2(targetY, targetX) - Math.atan2(-d4, Rp);
  const j1 = theta1 * 180 / Math.PI;

  if (j1 < jointsConfig[1].min || j1 > jointsConfig[1].max) {
    return { error: 'FUERA DE ALCANCE', msg: 'Rotacion base J1 fuera de limites (-165 a 165).' };
  }

  const xc = Rp - a2;
  const zc = targetZ - (d1 + d2);

  const psi_numerator = xc * xc + zc * zc - L1 * L1 - L2 * L2;
  const psi_denominator = 2.0 * L1 * L2;
  let psi = psi_numerator / psi_denominator;

  if (psi > 1.0 && psi <= 1.005) {
    psi = 1.0;
  } else if (psi < -1.0 && psi >= -1.005) {
    psi = -1.0;
  }

  if (psi < -1.0 || psi > 1.0) {
    return { error: 'FUERA DE ALCANCE', msg: 'Coordenadas fuera del alcance del brazo.' };
  }

  const solutions = [-1, 1];
  for (let s of solutions) {
    const sinTheta3_star = s * Math.sqrt(1.0 - psi * psi);
    const theta3_star = Math.atan2(sinTheta3_star, psi);

    const k1 = L1 + L2 * Math.cos(theta3_star);
    const k2 = L2 * Math.sin(theta3_star);
    const theta2_star = Math.atan2(k1 * zc - k2 * xc, k1 * xc + k2 * zc);

    const theta2 = theta2_star;
    const theta3 = theta3_star;

    const j2 = theta2 * 180 / Math.PI;
    const j3 = theta3 * 180 / Math.PI;

    if (j2 >= jointsConfig[2].min && j2 <= jointsConfig[2].max &&
        j3 >= jointsConfig[3].min && j3 <= jointsConfig[3].max) {
      return { j1, j2, j3 };
    }
  }

  return { error: 'FUERA DE ALCANCE', msg: 'Posicion cinematicamente inalcanzable por limites de articulacion.' };
}
\end{lstlisting}

\section{Código del Generador de Trayectorias de Simulación (trajectory\_generator.py)}
\label{ap:traj_py}
\begin{lstlisting}[language=Python,caption={\small Generación de perfiles S-curve por interpolación cosenoidal y control de keyframes (trajectory\_generator.py).},label={lst:traj_py}]
#!/usr/bin/env python3
import rclpy
from rclpy.node import Node
from std_msgs.msg import Float64MultiArray
import time
import math

class TrajectoryGenerator(Node):
    def __init__(self):
        """
        Inicializa el nodo de ROS 2, configura el publicador de comandos de control,
        define los keyframes de la secuencia de corte industrial y arranca el temporizador de interpolacion.
        """
        super().__init__('trajectory_generator')
        self.publisher_ = self.create_publisher(Float64MultiArray, '/arm_controller/commands', 10)

        self.rate = 50.0
        self.timer = self.create_timer(1.0 / self.rate, self.timer_callback)

        self.keyframes = [
            {"pose": [0.0, 0.0, 0.0], "time": 3.0},
            {"pose": [-0.6, 0.4, -0.4], "time": 3.0},
            {"pose": [0.6, 0.4, -0.4], "time": 4.0},
            {"pose": [0.6, 0.6, -0.7], "time": 2.0},
            {"pose": [-0.6, 0.6, -0.7], "time": 4.0},
            {"pose": [0.0, 0.0, 0.0], "time": 3.0}
        ]

        self.current_keyframe_idx = 0
        self.step_counter = 0
        self.total_steps_for_transition = 0
        self.start_pose = [0.0, 0.0, 0.0]
        self.target_pose = self.keyframes[0]["pose"]

        self.get_logger().info("Nodo de trayectoria inicializado. Iniciando ciclo de corte...")
        self.prepare_next_transition()

    def prepare_next_transition(self):
        """
        Prepara los parametros cinematicos y de interpolacion para la transicion hacia el
        siguiente keyframe de la secuencia, incrementando el indice de trayectoria.
        """
        self.start_pose = list(self.target_pose)
        keyframe = self.keyframes[self.current_keyframe_idx]
        self.target_pose = keyframe["pose"]
        duration = keyframe["time"]

        self.total_steps_for_transition = int(duration * self.rate)
        self.step_counter = 0

        self.get_logger().info(
            f"Transicionando a Keyframe {self.current_keyframe_idx}: {self.target_pose} en {duration}s"
        )

        self.current_keyframe_idx = (self.current_keyframe_idx + 1) % len(self.keyframes)

    def timer_callback(self):
        """
        Ejecuta periodicamente la interpolacion cosenoidal (S-curve) entre poses conjuntas
        sucesivas y publica el vector de comandos correspondiente a la tasa del bucle.
        """
        if self.step_counter >= self.total_steps_for_transition:
            self.prepare_next_transition()

        t = float(self.step_counter) / float(self.total_steps_for_transition)
        smooth_t = (1.0 - math.cos(t * math.pi)) / 2.0

        current_command = []
        for start, target in zip(self.start_pose, self.target_pose):
            val = start + (target - start) * smooth_t
            current_command.append(val)

        msg = Float64MultiArray()
        msg.data = current_command
        self.publisher_.publish(msg)

        self.step_counter += 1


def main(args=None):
    """
    Funcion de entrada que inicializa el entorno de ROS 2, instancia el generador de trayectoria,
    mantiene el ciclo de ejecucion activo y maneja el apagado limpio del nodo.
    """
    rclpy.init(args=args)
    node = TrajectoryGenerator()
    try:
        rclpy.spin(node)
    except KeyboardInterrupt:
        node.get_logger().info("Generador de trayectoria detenido por el usuario.")
    finally:
        node.destroy_node()
        rclpy.shutdown()

if __name__ == '__main__':
    main()
\end{lstlisting}

\section{Código del Puente de Comunicación ROS 2 - Interfaz Web (web\_bridge.py)}
\label{ap:bridge_py}
\begin{lstlisting}[language=Python,caption={\small Conversión de unidades y envío al simulador (web\_bridge.py).},label={lst:bridge_py}]
def publish_joints(self, j1_deg, j2_deg, j3_deg):
    j1_rad = float(j1_deg) * math.pi / 180.0
    j2_rad = float(j2_deg) * math.pi / 180.0
    j3_rad = float(j3_deg) * math.pi / 180.0
    msg = Float64MultiArray()
    msg.data = [j1_rad, j2_rad, j3_rad]
    self.publisher_.publish(msg)
    self.get_logger().info(f'Publicado en ROS -> J1: {j1_rad:.4f} rad, J2: {j2_rad:.4f} rad, J3: {j3_rad:.4f} rad')
\end{lstlisting}

\section{Código del Lanzamiento de Simulación en ROS 2 (gazebo.launch.py)}
\label{ap:launch_py}
\begin{lstlisting}[language=Python,caption={\small Script de lanzamiento y control secuencial en ROS 2 (gazebo.launch.py).},label={lst:launch_py}]
import os
from ament_index_python.packages import get_package_share_directory
from launch import LaunchDescription
from launch.actions import IncludeLaunchDescription, RegisterEventHandler
from launch.event_handlers import OnProcessExit
from launch.launch_description_sources import PythonLaunchDescriptionSource
from launch_ros.actions import Node
import xacro

def generate_launch_description():
    pkg_share = get_package_share_directory('arm_simulation')
    xacro_file = os.path.join(pkg_share, 'urdf', 'arm.urdf.xacro')

    pkg_share_parent = os.path.dirname(pkg_share)
    default_gazebo_models = '/usr/share/gazebo-11/models:/usr/share/gazebo/models'
    if 'GAZEBO_MODEL_PATH' in os.environ:
        os.environ['GAZEBO_MODEL_PATH'] = pkg_share_parent + ':' + default_gazebo_models + ':' + os.environ['GAZEBO_MODEL_PATH']
    else:
        os.environ['GAZEBO_MODEL_PATH'] = pkg_share_parent + ':' + default_gazebo_models

    os.environ['GAZEBO_MODEL_DATABASE_URI'] = ''

    import re
    robot_description_config = xacro.process_file(xacro_file)
    robot_desc = robot_description_config.toxml()
    robot_desc = re.sub(r'<!--.*?-->', '', robot_desc, flags=re.DOTALL)

    robot_state_publisher_node = Node(
        package='robot_state_publisher',
        executable='robot_state_publisher',
        name='robot_state_publisher',
        output='screen',
        parameters=[{'robot_description': robot_desc, 'use_sim_time': True}]
    )

    gazebo = IncludeLaunchDescription(
        PythonLaunchDescriptionSource([os.path.join(
            get_package_share_directory('gazebo_ros'), 'launch', 'gazebo.launch.py')]),
        launch_arguments={'verbose': 'true'}.items()
    )

    spawn_entity = Node(
        package='gazebo_ros',
        executable='spawn_entity.py',
        arguments=['-topic', 'robot_description', '-entity', 'arm_3gdl', '-timeout', '120'],
        output='screen'
    )

    spawn_joint_state_broadcaster = Node(
        package="controller_manager",
        executable="spawner",
        arguments=["joint_state_broadcaster", "--controller-manager", "/controller_manager"],
        output='screen'
    )

    spawn_arm_controller = Node(
        package="controller_manager",
        executable="spawner",
        arguments=["arm_controller", "--controller-manager", "/controller_manager"],
        output='screen'
    )

    return LaunchDescription([
        gazebo,
        robot_state_publisher_node,
        spawn_entity,
        RegisterEventHandler(
            event_handler=OnProcessExit(
                target_action=spawn_entity,
                on_exit=[spawn_joint_state_broadcaster],
            )
        ),
        RegisterEventHandler(
            event_handler=OnProcessExit(
                target_action=spawn_joint_state_broadcaster,
                on_exit=[spawn_arm_controller],
            )
        )
    ])
\end{lstlisting}

\section{Archivo de Descripción Estructural y Actuadores Xacro (arm.urdf.xacro)}
\label{ap:urdf_xacro}
\begin{lstlisting}[language=XML,caption={\small Modelado estructural, cinemático y configuración de ros2\_control en Xacro (arm.urdf.xacro).},label={lst:urdf_xacro}]
<?xml version="1.0"?>
<robot xmlns:xacro="http://www.ros.org/wiki/xacro" name="arm_3gdl">

  <!-- ==================== Materiales de Visualizacion ==================== -->
  <material name="abb_orange">
    <color rgba="1.0 0.45 0.0 1.0"/>
  </material>
  <material name="abb_black">
    <color rgba="0.12 0.12 0.12 1.0"/>
  </material>
  <material name="silver">
    <color rgba="0.75 0.75 0.75 1.0"/>
  </material>
  <material name="graphite">
    <color rgba="0.25 0.25 0.25 1.0"/>
  </material>
  <material name="brass">
    <color rgba="0.8 0.6 0.2 1.0"/>
  </material>
  <material name="laser_red">
    <color rgba="1.0 0.0 0.0 0.8"/>
  </material>

  <!-- ==================== Base Fija (Mundo) ==================== -->
  <link name="world"/>
  <joint name="world_to_base" type="fixed">
    <parent link="world"/>
    <child link="base_link"/>
    <origin xyz="0 0 0" rpy="1.57079632679 0 0"/>
  </joint>

  <!-- ==================== Eslabon Base ==================== -->
  <link name="base_link">
    <inertial>
      <mass value="1.0"/>
      <origin xyz="0 0 0" rpy="0 0 0"/>
      <inertia ixx="0.005" ixy="0.0" ixz="0.0" iyy="0.005" iyz="0.0" izz="0.005"/>
    </inertial>
    <visual>
      <origin xyz="-0.691316 -0.869483 -1.350860" rpy="0 0 0"/>
      <geometry>
        <mesh filename="package://arm_simulation/meshes/custom_arm/base.stl" scale="0.001 0.001 0.001"/>
      </geometry>
      <material name="abb_orange"/>
    </visual>
  </link>

  <!-- ==================== Articulacion 1 (Rotacion de Base) ==================== -->
  <joint name="joint1" type="revolute">
    <origin rpy="-1.57079632679 0 0" xyz="0 0.060 0"/>
    <parent link="base_link"/>
    <child link="link1"/>
    <axis xyz="0 0 1"/>
    <limit lower="-2.87979" upper="2.87979" effort="20.0" velocity="1.0"/>
    <dynamics damping="1.0" friction="1.0"/>
  </joint>

  <!-- ==================== Eslabon 1 ==================== -->
  <link name="link1">
    <inertial>
      <mass value="0.5"/>
      <origin xyz="0 0 0" rpy="0 0 0"/>
      <inertia ixx="0.002" ixy="0.0" ixz="0.0" iyy="0.002" iyz="0.0" izz="0.002"/>
    </inertial>
    <visual>
      <origin xyz="-0.691316 1.350860 -0.929483" rpy="1.57079632679 0 0"/>
      <geometry>
        <mesh filename="package://arm_simulation/meshes/custom_arm/link1.stl" scale="0.001 0.001 0.001"/>
      </geometry>
      <material name="abb_orange"/>
    </visual>
  </link>

  <!-- ==================== Articulacion 2 (Hombro) ==================== -->
  <joint name="joint2" type="revolute">
    <origin rpy="1.57079632679 0 0" xyz="0.014915 0 0.036173"/>
    <parent link="link1"/>
    <child link="link2"/>
    <axis xyz="0 0 1"/>
    <limit lower="-1.91986" upper="1.91986" effort="25.0" velocity="1.0"/>
    <dynamics damping="1.0" friction="1.0"/>
  </joint>

  <!-- ==================== Eslabon 2 ==================== -->
  <link name="link2">
    <inertial>
      <mass value="0.5"/>
      <origin xyz="0 0 0" rpy="0 0 0"/>
      <inertia ixx="0.002" ixy="0.0" ixz="0.0" iyy="0.002" iyz="0.0" izz="0.002"/>
    </inertial>
    <visual>
      <origin xyz="-0.706231 -0.965655 -1.350860" rpy="0 0 0"/>
      <geometry>
        <mesh filename="package://arm_simulation/meshes/custom_arm/link2.stl" scale="0.001 0.001 0.001"/>
      </geometry>
      <material name="abb_orange"/>
    </visual>
  </link>

  <!-- ==================== Articulacion 3 (Codo) ==================== -->
  <joint name="joint3" type="revolute">
    <origin rpy="0 0 0" xyz="0.075256 0.125332 0"/>
    <parent link="link2"/>
    <child link="link3"/>
    <axis xyz="0 0 1"/>
    <limit lower="-1.91986" upper="1.22173" effort="20.0" velocity="1.5"/>
    <dynamics damping="1.0" friction="1.0"/>
  </joint>

  <!-- ==================== Eslabon 3 ==================== -->
  <link name="link3">
    <inertial>
      <mass value="0.3"/>
      <origin xyz="0 0 0" rpy="0 0 0"/>
      <inertia ixx="0.001" ixy="0.0" ixz="0.0" iyy="0.001" iyz="0.0" izz="0.001"/>
    </inertial>
    <visual>
      <origin xyz="-0.781487 -1.090987 -1.350860" rpy="0 0 0"/>
      <geometry>
        <mesh filename="package://arm_simulation/meshes/custom_arm/link3.stl" scale="0.001 0.001 0.001"/>
      </geometry>
      <material name="abb_orange"/>
    </visual>
  </link>

  <!-- ==================== Acoplamiento de Herramienta (Flange y Tool0) ==================== -->
  <link name="flange"/>
  <joint type="fixed" name="joint_3-flange">
    <origin xyz="0.158 -0.030 -0.004" rpy="0 0 -0.191392"/>
    <parent link="link3"/>
    <child link="flange"/>
  </joint>

  <link name="tool0"/>
  <joint name="flange-tool0" type="fixed">
    <origin xyz="0 0 0" rpy="0 1.57079632679 0" />
    <parent link="flange" />
    <child link="tool0" />
  </joint>

  <!-- ==================== Controladores y Plugins de Gazebo ==================== -->
  <gazebo>
    <plugin name="gazebo_ros2_control" filename="libgazebo_ros2_control.so">
      <parameters>$(find arm_simulation)/config/controllers.yaml</parameters>
    </plugin>
  </gazebo>

  <ros2_control name="GazeboSystem" type="system">
    <hardware>
      <plugin>gazebo_ros2_control/GazeboSystem</plugin>
    </hardware>
    <joint name="joint1">
      <command_interface name="position">
        <param name="min">-2.87979</param>
        <param name="max">2.87979</param>
      </command_interface>
      <state_interface name="position">
        <param name="initial_value">0.0</param>
      </state_interface>
      <state_interface name="velocity"/>
    </joint>
    <joint name="joint2">
      <command_interface name="position">
        <param name="min">-1.91986</param>
        <param name="max">1.91986</param>
      </command_interface>
      <state_interface name="position">
        <param name="initial_value">0.0</param>
      </state_interface>
      <state_interface name="velocity"/>
    </joint>
    <joint name="joint3">
      <command_interface name="position">
        <param name="min">-1.91986</param>
        <param name="max">1.22173</param>
      </command_interface>
      <state_interface name="position">
        <param name="initial_value">0.0</param>
      </state_interface>
      <state_interface name="velocity"/>
    </joint>
  </ros2_control>

  <!-- ==================== Macro para Sistemas de Referencia (RGB) ==================== -->
  <xacro:macro name="visualize_frame" params="parent prefix xyz:='0 0 0' rpy:='0 0 0' scale:=0.05">
    <!-- Eje X (Rojo) -->
    <link name="${prefix}_frame_x">
      <inertial>
        <mass value="0.0001"/>
        <origin xyz="0 0 0" rpy="0 0 0"/>
        <inertia ixx="1e-9" ixy="0" ixz="0" iyy="1e-9" iyz="0" izz="1e-9"/>
      </inertial>
      <visual>
        <origin xyz="${scale/4} 0 0" rpy="0 1.57079632679 0"/>
        <geometry>
          <cylinder radius="${scale*0.01}" length="${scale/2}"/>
        </geometry>
        <material name="red">
          <color rgba="1.0 0.0 0.0 0.8"/>
        </material>
      </visual>
    </link>
    <joint name="${prefix}_frame_x_joint" type="fixed">
      <origin xyz="${xyz}" rpy="${rpy}"/>
      <parent link="${parent}"/>
      <child link="${prefix}_frame_x"/>
    </joint>
    <gazebo reference="${prefix}_frame_x">
      <material>Gazebo/Red</material>
    </gazebo>

    <!-- Eje Y (Verde) -->
    <link name="${prefix}_frame_y">
      <inertial>
        <mass value="0.0001"/>
        <origin xyz="0 0 0" rpy="0 0 0"/>
        <inertia ixx="1e-9" ixy="0" ixz="0" iyy="1e-9" iyz="0" izz="1e-9"/>
      </inertial>
      <visual>
        <origin xyz="0 ${scale/4} 0" rpy="-1.57079632679 0 0"/>
        <geometry>
          <cylinder radius="${scale*0.01}" length="${scale/2}"/>
        </geometry>
        <material name="green">
          <color rgba="0.0 1.0 0.0 0.8"/>
        </material>
      </visual>
    </link>
    <joint name="${prefix}_frame_y_joint" type="fixed">
      <origin xyz="${xyz}" rpy="${rpy}"/>
      <parent link="${parent}"/>
      <child link="${prefix}_frame_y"/>
    </joint>
    <gazebo reference="${prefix}_frame_y">
      <material>Gazebo/Green</material>
    </gazebo>

    <!-- Eje Z (Azul) -->
    <link name="${prefix}_frame_z">
      <inertial>
        <mass value="0.0001"/>
        <origin xyz="0 0 0" rpy="0 0 0"/>
        <inertia ixx="1e-9" ixy="0" ixz="0" iyy="1e-9" iyz="0" izz="1e-9"/>
      </inertial>
      <visual>
        <origin xyz="0 0 ${scale/4}" rpy="0 0 0"/>
        <geometry>
          <cylinder radius="${scale*0.01}" length="${scale/2}"/>
        </geometry>
        <material name="blue">
          <color rgba="0.0 0.0 1.0 0.8"/>
        </material>
      </visual>
    </link>
    <joint name="${prefix}_frame_z_joint" type="fixed">
      <origin xyz="${xyz}" rpy="${rpy}"/>
      <parent link="${parent}"/>
      <child link="${prefix}_frame_z"/>
    </joint>
    <gazebo reference="${prefix}_frame_z">
      <material>Gazebo/Blue</material>
    </gazebo>
  </xacro:macro>

  <!-- Instanciacion de los Sistemas de Referencia en cada Articulacion (D-H Secuencial) -->
  <xacro:visualize_frame parent="base_link" prefix="system0" xyz="0 0 0" rpy="-1.57079632679 0 0" scale="0.08"/>
  <xacro:visualize_frame parent="link1" prefix="system1" xyz="0 0 0" rpy="0 0 0" scale="0.06"/>
  <xacro:visualize_frame parent="link2" prefix="system2" xyz="0 0 0" rpy="0 0 1.030041" scale="0.06"/>
  <xacro:visualize_frame parent="link3" prefix="system3" xyz="0 0 0" rpy="0 0 -0.191392" scale="0.06"/>
  <xacro:visualize_frame parent="flange" prefix="system4" xyz="0 0 0" rpy="0 0 0" scale="0.06"/>

  <!-- ==================== Camara de Monitoreo (Gazebo) ==================== -->
  <link name="camera_link">
    <inertial>
      <mass value="0.1"/>
      <origin xyz="0 0 0" rpy="0 0 0"/>
      <inertia ixx="1e-4" ixy="0" ixz="0" iyy="1e-4" iyz="0" izz="1e-4"/>
    </inertial>
    <visual>
      <!-- Cuerpo de la camara -->
      <origin xyz="0 0 0" rpy="0 0 0"/>
      <geometry>
        <box size="0.05 0.05 0.05"/>
      </geometry>
      <material name="graphite"/>
    </visual>
    <visual>
      <!-- Lente de la camara -->
      <origin xyz="0.03 0 0" rpy="0 1.57079632679 0"/>
      <geometry>
        <cylinder radius="0.015" length="0.02"/>
      </geometry>
      <material name="abb_black"/>
    </visual>
  </link>

  <joint name="camera_joint" type="fixed">
    <!-- Ubicada sobre el eje X positivo de system0, conectada a base_link para evitar fallos de spawn en Gazebo -->
    <origin xyz="0.11 0.025 0.0" rpy="-1.57079632679 0 0"/>
    <parent link="base_link"/>
    <child link="camera_link"/>
  </joint>

  <!-- Plugin de Gazebo para transmitir la imagen en ROS 2 -->
  <gazebo reference="camera_link">
    <sensor type="camera" name="work_camera">
      <update_rate>15.0</update_rate>
      <visualize>true</visualize>
      <camera name="head">
        <horizontal_fov>1.085595</horizontal_fov>
        <image>
          <width>640</width>
          <height>480</height>
          <format>R8G8B8</format>
        </image>
        <clip>
          <near>0.05</near>
          <far>8.0</far>
        </clip>
      </camera>
      <plugin name="camera_controller" filename="libgazebo_ros_camera.so">
        <frame_name>camera_link</frame_name>
      </plugin>
    </sensor>
  </gazebo>

  <!-- Sistemas de referencia de la camara (Base y Lente) -->
  <xacro:visualize_frame parent="base_link" prefix="camera_base" xyz="0.11 0.0 0.0" rpy="-1.57079632679 0 0" scale="0.04"/>
  <xacro:visualize_frame parent="base_link" prefix="camera_lens" xyz="0.15 0.025 0.0" rpy="-1.57079632679 0 0" scale="0.02"/>

</robot>
\end{lstlisting}

\section{Script de Automatizacion de Configuracion y Compilacion (setup.sh)}
\label{ap:setup_sh}
\begin{lstlisting}[language=bash,caption={\small Script de inicialización y compilación del entorno (setup.sh).},label={lst:setup_sh}]
#!/bin/bash
set -e
GREEN='\033[0;32m'
BLUE='\033[0;34m'
YELLOW='\033[1;33m'
RED='\033[0;31m'
NC='\033[0m'

echo -e "${BLUE}===============================================${NC}"
echo -e "${GREEN}      KiraOne - SETUP & COMPILATION SYSTEM     ${NC}"
echo -e "${BLUE}===============================================${NC}"

if ! command -v distrobox &> /dev/null; then
    echo -e "${RED}Error: distrobox no esta instalado.${NC}"
    exit 1
fi

echo -e "${BLUE}[1/4] Verificando contenedor Distrobox...${NC}"
if ! distrobox list | grep -q "ros2-humble"; then
    distrobox create -n ros2-humble -i ubuntu:22.04 -y
fi
echo -e "${GREEN}[OK] Contenedor 'ros2-humble' detectado.${NC}"

echo -e "${BLUE}[2/4] Compilando Workspace de ROS2...${NC}"
SCRIPT_DIR="$( cd "$( dirname "${BASH_SOURCE[0]}" )" &> /dev/null && pwd )"

distrobox enter ros2-humble -- bash -c "
  cd '$SCRIPT_DIR/ros2_ws' && \
  source /opt/ros/humble/setup.bash && \
  colcon build --packages-select arm_simulation
"
echo -e "${GREEN}[OK] Workspace compilado exitosamente.${NC}"

echo -e "${BLUE}[3/4] Configurando permisos de ejecucion...${NC}"
chmod +x "$SCRIPT_DIR/ros2_ws/src/arm_simulation/scripts/trajectory_generator.py"
echo -e "${GREEN}[OK] Permisos configurados.${NC}"

echo -e "${BLUE}[4/4] Generando scripts de ejecucion rapida...${NC}"
# Generacion automatica de launch_gazebo.sh, run_trajectory.sh y start_ui.sh ...
\end{lstlisting}

\section{Lanzadores Rápidos de la Simulación y de la Interfaz Web}
\label{ap:runners_sh}
\begin{lstlisting}[language=bash,caption={\small Script lanzador de Gazebo Classic (launch\_gazebo.sh).},label={lst:runners_sh}]
#!/bin/bash
SCRIPT_DIR="$( cd "$( dirname "${BASH_SOURCE[0]}" )" &> /dev/null && pwd )"
echo "Iniciando Gazebo Classic (KiraOne)..."
distrobox enter ros2-humble -- bash -c "
  export PATH=\$(echo \$PATH | tr ':' '\n' | grep -v miniconda | tr '\n' ':')
  source /opt/ros/humble/setup.bash
  source '$SCRIPT_DIR/ros2_ws/install/setup.bash'
  ros2 launch arm_simulation gazebo.launch.py
"
\end{lstlisting}

\begin{lstlisting}[language=bash,caption={\small Script ejecutor de la trayectoria cosenoidal (run\_trajectory.sh).}]
#!/bin/bash
SCRIPT_DIR="$( cd "$( dirname "${BASH_SOURCE[0]}" )" &> /dev/null && pwd )"
echo "Ejecutando generador de trayectoria (KiraOne)..."
distrobox enter ros2-humble -- bash -c "
  export PATH=\$(echo \$PATH | tr ':' '\n' | grep -v miniconda | tr '\n' ':')
  source /opt/ros/humble/setup.bash
  source '$SCRIPT_DIR/ros2_ws/install/setup.bash'
  /usr/bin/python3 '$SCRIPT_DIR/ros2_ws/src/arm_simulation/scripts/trajectory_generator.py'
"
\end{lstlisting}

\begin{lstlisting}[language=bash,caption={\small Script de inicio del servidor y de la interfaz de control (start\_ui.sh).}]
#!/bin/bash
SCRIPT_DIR="$( cd "$( dirname "${BASH_SOURCE[0]}" )" &> /dev/null && pwd )"
echo "Iniciando servidor web y puente ROS 2 para KiraOne en http://localhost:8080..."

if command -v xdg-open &> /dev/null; then
  (sleep 1.5 && xdg-open "http://localhost:8080") &
elif command -v sensible-browser &> /dev/null; then
  (sleep 1.5 && sensible-browser "http://localhost:8080") &
fi

distrobox enter ros2-humble -- bash -c "
  export PATH=\$(echo \$PATH | tr ':' '\n' | grep -v miniconda | tr '\n' ':')
  source /opt/ros/humble/setup.bash
  source '$SCRIPT_DIR/ros2_ws/install/setup.bash'
  python3 '$SCRIPT_DIR/ui/web_bridge.py'
"
\end{lstlisting}

\section{Código del Firmware de Control de Bajo Nivel y Láser en el ESP32 (main.cpp)}
\label{ap:main_cpp}
\begin{lstlisting}[language=C++,caption={\small Firmware embebido en C++ para control de servomotores y conmutación láser en el ESP32 (main.cpp).},label={lst:main_cpp}]
#include <Arduino.h>
#include <ESP32Servo.h>
#include <math.h>

#define SERVO_BASE_PIN     27
#define SERVO_SHOULDER_PIN 14
#define SERVO_ELBOW_PIN    21
#define LASER_PIN          22

// Constantes fisicas del brazo (metros)
const float L0 = 0.150f;
const float L1 = 0.200f;
const float L2 = 0.200f;

// Limites de las articulaciones (radianes)
const float MIN_THETA1 = -1.5708f; // -90 grados
const float MAX_THETA1 =  1.5708f; // +90 grados
const float MIN_THETA2 = -0.1745f; // -10 grados
const float MAX_THETA2 =  1.5708f; // +90 grados
const float MIN_THETA3 = -2.0944f; // -120 grados
const float MAX_THETA3 =  0.1745f; // +10 grados

const unsigned long SERIAL_TIMEOUT_MS = 1000;

Servo servoBase;
Servo servoShoulder;
Servo servoElbow;

// Variables de estado (microsegundos correspondientes a la posicion del servo)
int currentBaseUs = 1500;
int currentShoulderUs = 1500;
int currentElbowUs = 1500;

unsigned long lastPacketTime = 0;
bool isEmergency = false;

/*
 * Convierte un angulo articular en radianes al correspondiente ancho de pulso PWM en microsegundos,
 * limitando el resultado entre los margenes fisicos y de seguridad del servomotor.
 */
int angleToMicroseconds(float angle_rad, float min_rad, float max_rad) {
  float clamped = constrain(angle_rad, min_rad, max_rad);
  float us = 1500.0f + clamped * (2000.0f / PI);
  return constrain((int)us, 500, 2500);
}

/*
 * Resuelve analiticamente el modelo cinematico inverso para un brazo de 3 GDL bajo
 * la configuracion geometrica de codo arriba a partir de una posicion cartesiana (x, y, z).
 */
bool solveInverseKinematics(float x, float y, float z, float &t1, float &t2, float &t3) {
  t1 = atan2(y, x);

  float r = sqrt(x * x + y * y);
  float z_prime = z - L0;

  float psi_numerator = x * x + y * y + z_prime * z_prime - L1 * L1 - L2 * L2;
  float psi_denominator = 2.0f * L1 * L2;
  float psi = psi_numerator / psi_denominator;

  if (psi < -1.0f || psi > 1.0f) {
    return false;
  }

  t3 = atan2(-sqrt(1.0f - psi * psi), psi);

  float k1 = L1 + L2 * cos(t3);
  float k2 = L2 * sin(t3);
  t2 = atan2(k1 * z_prime - k2 * r, k1 * r + k2 * z_prime);

  return true;
}

/*
 * Activa de forma inmediata el estado de emergencia, apaga el emisor laser de corte
 * y desacopla la senal de los servomotores para liberar la tension de la estructura.
 */
void triggerEmergency(const char* reason) {
  isEmergency = true;
  digitalWrite(LASER_PIN, LOW);
  
  servoBase.detach();
  servoShoulder.detach();
  servoElbow.detach();
  
  Serial.print("\n=== PARADA DE EMERGENCIA ACTIVA ===\nMotivo: ");
  Serial.println(reason);
  Serial.println("Reinicie el ESP32 para reestablecer el sistema.");
}

/*
 * Mueve todos los servos gradualmente en pasos simultaneos para evitar picos de corriente
 * bruscos que provoquen el apagado o proteccion de la fuente de alimentacion ATX.
 */
void moveToTarget(int targetBaseUs, int targetShoulderUs, int targetElbowUs, int stepUs = 15, int stepDelay = 15) {
  bool baseMoving = true;
  bool shoulderMoving = true;
  bool elbowMoving = true;
  
  while (baseMoving || shoulderMoving || elbowMoving) {
    if (isEmergency) {
      return;
    }
    
    // Alimentar el temporizador del watchdog de seguridad durante el recorrido
    lastPacketTime = millis();

    baseMoving = (currentBaseUs != targetBaseUs);
    shoulderMoving = (currentShoulderUs != targetShoulderUs);
    elbowMoving = (currentElbowUs != targetElbowUs);
    
    if (baseMoving) {
      if (abs(targetBaseUs - currentBaseUs) <= stepUs) {
        currentBaseUs = targetBaseUs;
      } else {
        currentBaseUs += (targetBaseUs > currentBaseUs) ? stepUs : -stepUs;
      }
      servoBase.writeMicroseconds(currentBaseUs);
    }
    
    if (shoulderMoving) {
      if (abs(targetShoulderUs - currentShoulderUs) <= stepUs) {
        currentShoulderUs = targetShoulderUs;
      } else {
        currentShoulderUs += (targetShoulderUs > currentShoulderUs) ? stepUs : -stepUs;
      }
      servoShoulder.writeMicroseconds(currentShoulderUs);
    }
    
    if (elbowMoving) {
      if (abs(targetElbowUs - currentElbowUs) <= stepUs) {
        currentElbowUs = targetElbowUs;
      } else {
        currentElbowUs += (targetElbowUs > currentElbowUs) ? stepUs : -stepUs;
      }
      servoElbow.writeMicroseconds(currentElbowUs);
    }
    
    if (baseMoving || shoulderMoving || elbowMoving) {
      delay(stepDelay);
    }
  }
}

void setup() {
  Serial.begin(115200);
  delay(500);
  Serial.println("\n=== ESP32 3-GDL Laser Arm Controller (Gradual Move) ===");

  pinMode(LASER_PIN, OUTPUT);
  digitalWrite(LASER_PIN, LOW);

  ESP32PWM::allocateTimer(0);
  ESP32PWM::allocateTimer(1);
  ESP32PWM::allocateTimer(2);
  
  servoBase.setPeriodHertz(50);
  servoShoulder.setPeriodHertz(50);
  servoElbow.setPeriodHertz(50);

  // Inicializar servos y asociar pines
  servoBase.attach(SERVO_BASE_PIN, 500, 2500);
  servoShoulder.attach(SERVO_SHOULDER_PIN, 500, 2500);
  servoElbow.attach(SERVO_ELBOW_PIN, 500, 2500);

  // Establecer posicion central por defecto al encender
  servoBase.writeMicroseconds(currentBaseUs);
  servoShoulder.writeMicroseconds(currentShoulderUs);
  servoElbow.writeMicroseconds(currentElbowUs);

  lastPacketTime = millis();
  Serial.println("Sistema listo. Esperando comandos seriales...");
}

void loop() {
  if (isEmergency) {
    digitalWrite(LASER_PIN, LOW);
    delay(100);
    return;
  }

  // Watchdog de seguridad: apaga el laser si se pierde la comunicacion serial
  if (millis() - lastPacketTime > SERIAL_TIMEOUT_MS) {
    digitalWrite(LASER_PIN, LOW);
  }

  if (Serial.available() > 0) {
    String line = Serial.readStringUntil('\n');
    line.trim();

    if (line.length() > 0) {
      // Comando manual para parada de emergencia
      if (line.equalsIgnoreCase("ESTOP") || line.equalsIgnoreCase("EMERGENCY")) {
        triggerEmergency("Comando manual recibido");
        return;
      }

      int idxX = line.indexOf("X:");
      int idxY = line.indexOf("Y:");
      int idxZ = line.indexOf("Z:");

      if (idxX != -1 && idxY != -1 && idxZ != -1) {
        int endX = line.indexOf(',', idxX);
        int endY = line.indexOf(',', idxY);
        int endZ = line.length();

        String xStr = line.substring(idxX + 2, endX != -1 ? endX : line.length());
        String yStr = line.substring(idxY + 2, endY != -1 ? endY : line.length());
        String zStr = line.substring(idxZ + 2, endZ);

        float raw_x = xStr.toFloat();
        float raw_y = yStr.toFloat();
        float raw_z = zStr.toFloat();

        float scale = 1.0f;
        if (abs(raw_x) > 2.0f || abs(raw_y) > 2.0f || abs(raw_z) > 2.0f) {
          scale = 0.001f;
        }

        float x = raw_x * scale;
        float y = raw_y * scale;
        float z = raw_z * scale;

        float theta1 = 0.0f, theta2 = 0.0f, theta3 = 0.0f;
        if (solveInverseKinematics(x, y, z, theta1, theta2, theta3)) {
          lastPacketTime = millis();
          digitalWrite(LASER_PIN, HIGH);

          int base_us = angleToMicroseconds(theta1, MIN_THETA1, MAX_THETA1);
          int shoulder_us = angleToMicroseconds(theta2, MIN_THETA2, MAX_THETA2);
          int elbow_us = angleToMicroseconds(theta3, MIN_THETA3, MAX_THETA3);

          // Movimiento de servos en pasos graduales suaves
          moveToTarget(base_us, shoulder_us, elbow_us);

          Serial.print("OK | Target: X=");
          Serial.print(x, 4);
          Serial.print(" Y=");
          Serial.print(y, 4);
          Serial.print(" Z=");
          Serial.print(z, 4);
          Serial.print(" | Joints: T1=");
          Serial.print(theta1 * 180.0f / PI, 1);
          Serial.print(" T2=");
          Serial.print(theta2 * 180.0f / PI, 1);
          Serial.print(" T3=");
          Serial.println(theta3 * 180.0f / PI, 1);
        } else {
          Serial.println("ERR: Posicion cartesiana fuera del area de alcance.");
        }
      } else {
        Serial.println("WARN: Formato invalido. Use 'X:val,Y:val,Z:val'");
      }
    }
  }
}
\end{lstlisting}

\section{Código del Algoritmo de Adelgazamiento Zhang-Suen (web\_bridge.py)}
\label{ap:zhang_py}
\begin{lstlisting}[language=Python,caption={\small Algoritmo de esqueletización vectorizada de Zhang-Suen (web\_bridge.py).},label={lst:zhang_py}]
def zhang_suen_thinning(img):
    import numpy as np
    im = (img > 0).astype(np.uint8)
    im = np.pad(im, 1, mode='constant', constant_values=0)
    
    while True:
        # Paso 1
        p2 = im[:-2, 1:-1]
        p3 = im[:-2, 2:]
        p4 = im[1:-1, 2:]
        p5 = im[2:, 2:]
        p6 = im[2:, 1:-1]
        p7 = im[2:, :-2]
        p8 = im[1:-1, :-2]
        p9 = im[:-2, :-2]
        
        B = p2 + p3 + p4 + p5 + p6 + p7 + p8 + p9
        
        A = ((p2 == 0) & (p3 == 1)).astype(np.uint8) + \
            ((p3 == 0) & (p4 == 1)).astype(np.uint8) + \
            ((p4 == 0) & (p5 == 1)).astype(np.uint8) + \
            ((p5 == 0) & (p6 == 1)).astype(np.uint8) + \
            ((p6 == 0) & (p7 == 1)).astype(np.uint8) + \
            ((p7 == 0) & (p8 == 1)).astype(np.uint8) + \
            ((p8 == 0) & (p9 == 1)).astype(np.uint8) + \
            ((p9 == 0) & (p2 == 1)).astype(np.uint8)
            
        cond1 = (B >= 2) & (B <= 6) & (A == 1) & (p2 * p4 * p6 == 0) & (p4 * p6 * p8 == 0)
        inner = im[1:-1, 1:-1]
        to_delete_1 = (inner == 1) & cond1
        
        im[1:-1, 1:-1] = np.where(to_delete_1, 0, inner)
        
        # Paso 2
        p2 = im[:-2, 1:-1]
        p3 = im[:-2, 2:]
        p4 = im[1:-1, 2:]
        p5 = im[2:, 2:]
        p6 = im[2:, 1:-1]
        p7 = im[2:, :-2]
        p8 = im[1:-1, :-2]
        p9 = im[:-2, :-2]
        
        B = p2 + p3 + p4 + p5 + p6 + p7 + p8 + p9
        A = ((p2 == 0) & (p3 == 1)).astype(np.uint8) + \
            ((p3 == 0) & (p4 == 1)).astype(np.uint8) + \
            ((p4 == 0) & (p5 == 1)).astype(np.uint8) + \
            ((p5 == 0) & (p6 == 1)).astype(np.uint8) + \
            ((p6 == 0) & (p7 == 1)).astype(np.uint8) + \
            ((p7 == 0) & (p8 == 1)).astype(np.uint8) + \
            ((p8 == 0) & (p9 == 1)).astype(np.uint8) + \
            ((p9 == 0) & (p2 == 1)).astype(np.uint8)
            
        cond2 = (B >= 2) & (B <= 6) & (A == 1) & (p2 * p4 * p8 == 0) & (p2 * p6 * p8 == 0)
        inner = im[1:-1, 1:-1]
        to_delete_2 = (inner == 1) & cond2
        
        im[1:-1, 1:-1] = np.where(to_delete_2, 0, inner)
        
        if not np.any(to_delete_1) and not np.any(to_delete_2):
            break
            
    return (im[1:-1, 1:-1] * 255).astype(np.uint8)
\end{lstlisting}

\section{Código de Verificación Cinemática y DH (test\_kinematics.py)}
\label{ap:test_kin}
\begin{lstlisting}[language=Python,caption={\small Comparación de modelos de cinemática directa y Denavit-Hartenberg (test\_kinematics.py).},label={lst:test_kin}]
import numpy as np

# Constantes del robot
d1 = 60.0
d2 = 36.173
a2 = 14.915
a3 = 146.190
a4 = 160.823
d4 = -4.0

phi2 = 1.030041
phi3 = -0.191392

def dh_matrix(theta, d, a, alpha):
    alpha_rad = np.radians(alpha) if isinstance(alpha, (int, float)) else alpha
    return np.array([
        [np.cos(theta), -np.sin(theta)*np.cos(alpha_rad),  np.sin(theta)*np.sin(alpha_rad), a*np.cos(theta)],
        [np.sin(theta),  np.cos(theta)*np.cos(alpha_rad), -np.cos(theta)*np.sin(alpha_rad), a*np.sin(theta)],
        [0,              np.sin(alpha_rad),               np.cos(alpha_rad),              d],
        [0,              0,                               0,                              1]
    ])

def fk_dh(theta1, theta2, theta3):
    theta2_star = theta2 + phi2
    theta3_star = theta3 + phi3 - phi2
    
    T1 = dh_matrix(theta1, d1 + d2, a2, 90)
    T2 = dh_matrix(theta2_star, 0, a3, 0)
    T3 = dh_matrix(theta3_star, d4, a4, 0)
    
    T = T1 @ T2 @ T3
    return T

def fk_analytical(theta1, theta2, theta3):
    theta2_star = theta2 + phi2
    theta3_star = theta3 + phi3 - phi2
    
    x_plane = a2 + a3 * np.cos(theta2_star) + a4 * np.cos(theta2_star + theta3_star)
    z_plane = d1 + d2 + a3 * np.sin(theta2_star) + a4 * np.sin(theta2_star + theta3_star)
    
    x = x_plane * np.cos(theta1) + d4 * np.sin(theta1)
    y = x_plane * np.sin(theta1) - d4 * np.cos(theta1)
    z = z_plane
    return np.array([x, y, z])

# Probar con angulos cero
t1, t2, t3 = 0.0, 0.0, 0.0
T_dh = fk_dh(t1, t2, t3)
pos_dh = T_dh[:3, 3]
pos_analyt = fk_analytical(t1, t2, t3)

print("D-H End Effector Position (3 matrices):")
print(pos_dh)
print("\nAnalytical End Effector Position:")
print(pos_analyt)
\end{lstlisting}

\section{Código de Validación de Alineación de Marcos TF (test\_tf\_alignment.py)}
\label{ap:test_tf}
\begin{lstlisting}[language=Python,caption={\small Validación de paralelismo de ejes longitudinales en base al árbol de transformadas TF (test\_tf\_alignment.py).},label={lst:test_tf}]
#!/usr/bin/env python3
import rclpy
from rclpy.node import Node
import tf2_ros
import math
import numpy as np

class TFAlignmentValidator(Node):
    def __init__(self):
        super().__init__('tf_alignment_validator')
        self.tf_buffer = tf2_ros.Buffer()
        self.tf_listener = tf2_ros.TransformListener(self.tf_buffer, self)
        
        # Esperar un poco a que se llene el buffer de TFs
        self.timer = self.create_timer(1.0, self.validate_alignment)
        self.get_logger().info("Validator node started. Waiting for TF data...")

    def validate_alignment(self):
        self.timer.cancel()
        
        frames = ['system1', 'system2', 'system3', 'flange']
        base_frame = 'link1'
        
        print("\n" + "="*70)
        print(" VALIDACION DE ALINEACION DE EJES X (VISTA LATERAL / PLANO X-Z de link1)")
        print("="*70)
        
        all_parallel = True
        
        for frame in frames:
            try:
                # Obtener la transformacion de link1 al sistema
                now = rclpy.time.Time()
                trans = self.tf_buffer.lookup_transform(base_frame, frame, now, timeout=rclpy.duration.Duration(seconds=2.0))
                
                # Obtener cuaternion
                q = trans.transform.rotation
                
                # Convertir cuaternion a matriz de rotacion
                R = self.quaternion_to_matrix(q.x, q.y, q.z, q.w)
                
                # El eje X del frame en coordenadas de link1 es la primera columna de R
                x_axis = R[:, 0]
                
                # Calcular angulo en el plano vertical X-Z de link1
                angle_rad = math.atan2(x_axis[2], x_axis[0])
                angle_deg = math.degrees(angle_rad)
                
                print(f"\nFrame: {frame}")
                print(f"  Vector eje X (en frame link1): [{x_axis[0]:.6f}, {x_axis[1]:.6f}, {x_axis[2]:.6f}]")
                print(f"  Desviacion lateral (Y): {x_axis[1]:.6f} (debe ser aprox 0)")
                print(f"  Angulo en plano X-Z (Pitch): {angle_deg:.4f} deg")
                
                if abs(x_axis[1]) > 0.005:
                    print("  [WARN] ADVERTENCIA: El eje X no esta en el mismo plano de movimiento (Y != 0)")
                    all_parallel = False
                elif abs(angle_deg) > 0.1:
                    print(f"  [WARN] ADVERTENCIA: El eje X no es paralelo al eje X horizontal (Angulo = {angle_deg:.4f} deg)")
                    all_parallel = False
                else:
                    print("  [OK] ALINEACION PERFECTA (eje X paralelo y horizontal)")
                    
            except Exception as e:
                self.get_logger().error(f"No se pudo obtener la transformacion para {frame}: {str(e)}")
                all_parallel = False
        
        print("\n" + "="*70)
        if all_parallel:
            print(" RESULTADO FINAL: [OK] TODOS LOS EJES X SON PARALELOS Y HORIZONTALES!")
            print(" El brazo esta perfectamente recto y alineado en la pose cero de la simulacion.")
        else:
            print(" RESULTADO FINAL: [ERR] EXISTEN DESALINEACIONES O DESVIACIONES FUERA DE LIMITES.")
        print("="*70 + "\n")
        
        rclpy.shutdown()

    def quaternion_to_matrix(self, x, y, z, w):
        return np.array([
            [1 - 2*(y**2 + z**2), 2*(x*y - w*z), 2*(x*z + w*y)],
            [2*(x*y + w*z), 1 - 2*(x**2 + z**2), 2*(y*z - w*x)],
            [2*(x*z - w*y), 2*(y*z + w*x), 1 - 2*(x**2 + y**2)]
        ])

def main(args=None):
    rclpy.init(args=args)
    node = TFAlignmentValidator()
    try:
        rclpy.spin(node)
    except SystemExit:
        pass

if __name__ == '__main__':
    main()
\end{lstlisting}

\section{Código de Verificación Analítica de URDF (test\_urdf\_tf\_alignment.py)}
\label{ap:test_urdf_tf}
\begin{lstlisting}[language=Python,caption={\small Verificación analítica de la rotación y coplanaridad de las articulaciones en el modelo URDF (test\_urdf\_tf\_alignment.py).},label={lst:test_urdf_tf}]
#!/usr/bin/env python3
import math
import numpy as np

def rpy_to_matrix(r, p, y):
    cr, sr = math.cos(r), math.sin(r)
    cp, sp = math.cos(p), math.sin(p)
    cy, sy = math.cos(y), math.sin(y)
    
    Rx = np.array([[1, 0, 0], [0, cr, -sr], [0, sr, cr]])
    Ry = np.array([[cp, 0, sp], [0, 1, 0], [-sp, 0, cp]])
    Rz = np.array([[cy, -sy, 0], [sy, cy, 0], [0, 0, 1]])
    
    return Rz @ Ry @ Rx

# Datos de origen de las articulaciones de la simulacion (URDF actual)
R_w_base = rpy_to_matrix(1.57079632679, 0, 0)
R_base_l1 = rpy_to_matrix(-1.57079632679, 0, 0)
R_l1_l2 = rpy_to_matrix(1.57079632679, 0, 0)
R_l2_l3 = rpy_to_matrix(0, 0, 0)
R_l3_flange = rpy_to_matrix(0, 0, 0)

# Marcos de referencia visuales (visualize_frame) en el URDF actual
R_base_sys0 = rpy_to_matrix(-1.57079632679, 0, 0)
R_l2_sys1 = rpy_to_matrix(0, 0, 0)
R_l3_sys2 = rpy_to_matrix(0, 0, 0)
R_flange_sys3 = rpy_to_matrix(0, 0, 0)

transforms = {
    "system1": R_l1_l2 @ R_l2_sys1,
    "system2": R_l1_l2 @ R_l2_l3 @ R_l3_sys2,
    "system3": R_l1_l2 @ R_l2_l3 @ R_l3_flange @ R_flange_sys3,
    "flange": R_l1_l2 @ R_l2_l3 @ R_l3_flange
}

print("\n" + "="*75)
print(" CALCULO ANALITICO DE ALINEACION DE EJES X (DESDE LA SIMULACION URDF)")
print(" Plano de referencia: vista lateral (plano X-Z de link1)")
print("="*75)

all_parallel = True

for name, R_l1_S in transforms.items():
    x_axis_in_l1 = R_l1_S[:, 0]
    y_deviation = x_axis_in_l1[1]
    angle_rad = math.atan2(x_axis_in_l1[2], x_axis_in_l1[0])
    angle_deg = math.degrees(angle_rad)
    
    print(f"\nSistema: {name}")
    print(f"  Direccion del eje X en link1: [{x_axis_in_l1[0]:.6f}, {x_axis_in_l1[1]:.6f}, {x_axis_in_l1[2]:.6f}]")
    print(f"  Desviacion lateral (Y): {y_deviation:.6f} (debe ser 0 para estar coplanar)")
    print(f"  Angulo lateral (Pitch): {angle_deg:.4f} deg (debe ser 0 deg para estar horizontal)")
    
    if abs(y_deviation) > 1e-5:
        print("  [ERR] ERROR: El eje X se sale del plano vertical (desviacion Y no es cero).")
        all_parallel = False
    elif abs(angle_deg) > 1e-5:
        print(f"  [ERR] ERROR: El eje X no es horizontal en la vista lateral (Angulo = {angle_deg:.4f} deg).")
        all_parallel = False
    else:
        print("  [OK] PERFECTO: El eje X es coplanar, paralelo y horizontal.")

print("\n" + "="*75)
if all_parallel:
    print(" VERIFICACION EXITOSA: [OK] TODOS LOS EJES X SON PARALELOS Y PERFECTAMENTE ALINEADOS!")
else:
    print(" VERIFICACION FALLIDA: [ERR] Existen desalineaciones entre los marcos.")
print("="*75 + "\n")
\end{lstlisting}

\section{Código de Lectura de Distancias de Marcos TF en Gazebo (get\_tf\_distances.py)}
\label{ap:get_tf_dist}
\begin{lstlisting}[language=Python,caption={\small Script de lectura y visualización de transformadas de articulaciones desde TF en Gazebo (get\_tf\_distances.py).},label={lst:get_tf_dist}]
#!/usr/bin/env python3
import rclpy
from rclpy.node import Node
import tf2_ros
import sys

class TFEchoAll(Node):
    def __init__(self):
        super().__init__('tf_echo_all')
        self.tf_buffer = tf2_ros.Buffer()
        self.tf_listener = tf2_ros.TransformListener(self.tf_buffer, self)
        self.timer = self.create_timer(1.0, self.query_all_tfs)
        self.run_once = False

    def query_all_tfs(self):
        if self.run_once:
            return
        
        frames = ['link1', 'link2', 'link3', 'flange']
        base = 'base_link'
        
        print("\n=== COORDENADAS DE LAS ARTICULACIONES DESDE TF EN GAZEBO ===")
        print(f"Sistema de referencia base: {base}\n")
        
        success = True
        for frame in frames:
            try:
                trans = self.tf_buffer.lookup_transform(base, frame, rclpy.time.Time())
                translation = trans.transform.translation
                rotation = trans.transform.rotation
                x, y, z, w = rotation.x, rotation.y, rotation.z, rotation.w
                
                print(f"-> {frame} respecto a {base}:")
                print(f"   Traslacion: [{translation.x:.6f}, {translation.y:.6f}, {translation.z:.6f}] m")
                print(f"               [{translation.x*1000:.3f}, {translation.y*1000:.3f}, {translation.z*1000:.3f}] mm")
                print(f"   Rotacion (Quaternion): xyzw = [{x:.6f}, {y:.6f}, {z:.6f}, {w:.6f}]\n")
                
            except Exception as e:
                print(f"Error buscando transformada de {base} a {frame}: {str(e)}")
                success = False
                
        if success:
            self.run_once = True
            sys.exit(0)

def main(args=None):
    rclpy.init(args=args)
    node = TFEchoAll()
    try:
        rclpy.spin(node)
    except SystemExit:
        pass
    except KeyboardInterrupt:
        pass
    finally:
        node.destroy_node()
        rclpy.shutdown()

if __name__ == '__main__':
    main()
\end{lstlisting}

\section{Código del Parser y Analizador de Parámetros Xacro (get\_gazebo\_distances.py)}
\label{ap:get_gz_dist}
\begin{lstlisting}[language=Python,caption={\small Script parser XML para extraer dimensiones y ángulos del archivo URDF/Xacro (get\_gazebo\_distances.py).},label={lst:get_gz_dist}]
import xml.etree.ElementTree as ET
import math
import os

def parse_urdf():
    urdf_path = "/home/reve/Documents/UC/robotica/finalProject/ros2_ws/src/arm_simulation/urdf/arm.urdf.xacro"
    if not os.path.exists(urdf_path):
        print(f"Error: No se encontro el archivo URDF en {urdf_path}")
        return

    tree = ET.parse(urdf_path)
    root = tree.getroot()

    print("=== DISTANCIAS Y LONGITUDES EXTRAIDAS DEL MODELO URDF EN GAZEBO ===\n")

    joint_world = root.find(".//joint[@name='world_to_base']")
    if joint_world is not None:
        origin = joint_world.find("origin")
        print(f"world_to_base:")
        print(f"  xyz: {origin.attrib.get('xyz')}")
        print(f"  rpy: {origin.attrib.get('rpy')}\n")

    joint1 = root.find(".//joint[@name='joint1']")
    if joint1 is not None:
        origin = joint1.find("origin")
        print(f"joint1 (Base a Link1):")
        print(f"  xyz: {origin.attrib.get('xyz')} (Desfase vertical base height d1 = 60 mm)")
        print(f"  rpy: {origin.attrib.get('rpy')}\n")

    joint2 = root.find(".//joint[@name='joint2']")
    if joint2 is not None:
        origin = joint2.find("origin")
        xyz = [float(x) for x in origin.attrib.get('xyz').split()]
        print(f"joint2 (Link1 a Hombro):")
        print(f"  xyz: {origin.attrib.get('xyz')} (Desfase radial a2 = {xyz[0]*1000:.3f} mm, Desfase vertical d2 = {xyz[2]*1000:.3f} mm)")
        print(f"  rpy: {origin.attrib.get('rpy')}\n")

    joint3 = root.find(".//joint[@name='joint3']")
    if joint3 is not None:
        origin = joint3.find("origin")
        xyz = [float(x) for x in origin.attrib.get('xyz').split()]
        length = math.sqrt(xyz[0]**2 + xyz[1]**2) * 1000
        angle = math.atan2(xyz[1], xyz[0]) * 180 / math.pi
        print(f"joint3 (Hombro a Codo):")
        print(f"  xyz: {origin.attrib.get('xyz')}")
        print(f"  Longitud del brazo (Link2) a3: {length:.3f} mm")
        print(f"  Angulo de inclinacion phi2: {angle:.6f} deg ({math.atan2(xyz[1], xyz[0]):.6f} rad)\n")

    joint_flange = root.find(".//joint[@name='joint_3-flange']")
    if joint_flange is not None:
        origin = joint_flange.find("origin")
        xyz = [float(x) for x in origin.attrib.get('xyz').split()]
        length = math.sqrt(xyz[0]**2 + xyz[1]**2) * 1000
        angle = math.atan2(xyz[1], xyz[0]) * 180 / math.pi
        print(f"joint_3-flange (Codo a Brida):")
        print(f"  xyz: {origin.attrib.get('xyz')}")
        print(f"  Longitud del antebrazo (Link3) a4: {length:.3f} mm")
        print(f"  Angulo de inclinacion phi3: {angle:.6f} deg ({math.atan2(xyz[1], xyz[0]):.6f} rad)")
        print(f"  Desfase transversal d4: {xyz[2]*1000:.3f} mm\n")

if __name__ == "__main__":
    parse_urdf()
\end{lstlisting}
\begin{thebibliography}{00}
\bibitem{b1} J. Craig, \textit{Introduction to Robotics: Mechanics and Control}, 4th ed. Pearson, 2017.
\bibitem{b2} I. Hinojosa, ``Repositorio de Código del Proyecto FP\_Robotics'', GitHub, 2026. Disponible en: \url{https://github.com/Reve7339/FP_Robotics}.
\end{thebibliography}

\end{document}
